% Formal-owned proof companion. Publication-quality mathematical exposition; not itself a publication or scientific-validation record.
\documentclass[11pt]{report}

\usepackage[T1]{fontenc}
\usepackage{amsmath,amssymb}
\usepackage{amsthm}
\usepackage{booktabs}
\usepackage[margin=1in]{geometry}
\usepackage[hidelinks]{hyperref}
\usepackage{microtype}
\usepackage{placeins}
\usepackage{seqsplit}
\usepackage{tabularx}

\newcolumntype{Y}{>{\raggedright\arraybackslash}X}
\newcolumntype{Z}{>{\raggedright\arraybackslash}p{0.23\textwidth}}
\def\UrlBreaks{\do\/\do-\do.\do_}
\Urlmuskip=0mu plus 1mu
\setlength{\emergencystretch}{3em}
\newcommand{\decl}[1]{\texttt{\seqsplit{#1}}}
\newcommand{\modulepath}[1]{\nolinkurl{#1}}

\newtheorem{theorem}{Theorem}[chapter]
\theoremstyle{definition}
\newtheorem{definition}[theorem]{Definition}
\newtheorem{openproblem}[theorem]{Open Problem}
\theoremstyle{remark}
\newtheorem{remark}[theorem]{Remark}

\title{Compatibility of Spectral and Operator Reductions\\
\large Admissible-Set Geometry, Alignment, Fourier Locality, and Exact Finite Decisions\\[0.5em]
\normalsize A Formal Proof Companion}
\author{Project Koios Formal}
\date{}

\begin{document}
\hypersetup{pageanchor=false}
\maketitle

\begin{abstract}
A spectral fit and an aligned operator fit need not select the same reduced Hamiltonian.
The central model-adequacy problem is therefore to determine, within each class of a
declared model hierarchy, whether the spectral and operator admissible sets contain one
common candidate; when they do not, their metric separation quantifies the obstruction.
This proof companion develops exact finite mathematics supporting that problem. It
distinguishes retained subspaces from retained coordinates, places alignment-family
selection inside the operator criterion, separates fixed-unitary Fourier equivariance
from pointwise-gauge mode mixing, and connects admissible-set geometry to exact and
interval-valued finite decisions. An exact three-channel calculus separates
reference-to-interpolant, interpolant-to-truncated, and reference-to-truncated residuals;
it proves their vector identity and pointwise, maximum, and root-mean-square triangle
bounds without asserting scalar-error additivity. An exact two-point rank-two
construction shows that a pointwise gauge need not commute with zero-mode Fourier
truncation, while finite-tail identities control only the deliberate-truncation leg.
The resulting framework supports common-witness compatibility and first-tested-model
selection without claiming that a first-principles calculation, runtime certificate, or
material validation has been completed. Declaration-level traceability, source identities, mechanization status,
and reproducibility metadata are collected in the appendices.
\end{abstract}

\clearpage
\hypersetup{pageanchor=true}
\pagenumbering{roman}
\tableofcontents
\clearpage
\pagenumbering{arabic}

\chapter{Spectral--Operator Compatibility as a Model-Adequacy Problem}
\label{ch:proof-architecture}

\section{The central question}

A reduced Hamiltonian may be judged by more than one criterion.  A spectral reduction
asks whether selected eigenvalues or derived band quantities are reproduced within
declared tolerances.  An operator reduction asks whether the reduced Hamiltonian
reproduces matrix information in a common aligned representation.  These requirements
are not equivalent.  Two models may have nearly identical selected eigenvalues while
differing in onsite terms, orbital couplings, hopping blocks, or their response to a
later perturbation.

The central model-adequacy question is therefore not whether two independently optimized
models are individually successful.  It is whether one candidate Hamiltonian satisfies
both requirements.

Let $J$ be a declared finite ordered inventory of model classes.  For each $j\in J$,
let $\mathfrak M_j$ be the corresponding candidate class and let $\theta$ denote a
candidate parameter in that class.  Write
\[
  H_j(\theta)\in\mathfrak M_j
\]
for the candidate Hamiltonian.  A spectral loss $\mathcal L_{E,j}(\theta)$ and an
aligned operator loss $\mathcal L_{H,j}(C,\theta)$ are evaluated against one declared
reference problem.  The map $C$ identifies candidate coordinates with reference
coordinates and must belong to a specified admissible alignment family
$\mathcal U_j$.

The two admissible sets are
\[
  \mathfrak A_{E,j}(\tau_E)
  =\left\{
      H_j(\theta)\in\mathfrak M_j:
      \mathcal L_{E,j}(\theta)\leq\tau_E
    \right\}
\]
and
\[
  \mathfrak A_{H,j}(\tau_H)
  =\left\{
      H_j(\theta)\in\mathfrak M_j:
      \inf_{C\in\mathcal U_j}
      \mathcal L_{H,j}(C,\theta)\leq\tau_H
    \right\}.
\]
When the alignment optimization is known to attain its infimum, the infimum may be
written as a minimum.  Attainment is a mathematical premise, not a consequence of
running an optimizer.

Exact compatibility within the model class is the proposition
\[
  \mathfrak A_{E,j}(\tau_E)
  \cap
  \mathfrak A_{H,j}(\tau_H)
  \neq\varnothing.
\]
When the alignment infimum is attained, this is equivalently the existence of one
parameter $\theta$ and one admissible alignment $C\in\mathcal U_j$ such that
\[
  \mathcal L_{E,j}(\theta)\leq\tau_E,
  \qquad
  \mathcal L_{H,j}(C,\theta)\leq\tau_H.
\]
The quantifiers are coupled.  A spectrally successful parameter, an operator-successful
parameter, and a favorable alignment chosen for some third parameter do not assemble
into a common witness.

This is the problem stated by the ICMSEP extended abstract that motivates the present
formal synthesis: determine whether spectral and operator reductions admit a common
Hamiltonian in each prescribed model class, quantify their separation when they do not,
and identify the first tested class that supports a common endpoint.  The present
companion develops exact finite mathematics supporting that question.  It does not
claim that the proposed bulk-silicon calculation has been executed or scientifically
validated.

\section{Separation and tolerance-level compatibility}

When the two admissible sets do not intersect, a metric on common canonical model
coordinates can measure the obstruction.  If $d_{\mathrm{model}}$ is that declared
metric, define
\[
  \delta_j^*
  =\inf_{\substack{
       H_E\in\mathfrak A_{E,j}(\tau_E)\\
       H_H\in\mathfrak A_{H,j}(\tau_H)}}
    d_{\mathrm{model}}(H_E,H_H).
\]
A positive certified lower bound on $\delta_j^*$ proves that the two sets are disjoint.
The converse direction requires care.  In an arbitrary metric space, distance zero does
not by itself produce an intersection point.  Exact equivalence between zero distance
and nonempty intersection requires suitable closure, compactness, attainment, or other
explicit hypotheses.  The finite geometric results used later state the assumptions
under which this converse is valid.

An application may also introduce a map tolerance $\tau_{\mathrm{map}}$ and declare
$\delta_j^*\leq\tau_{\mathrm{map}}$ to be an accepted approximate compatibility
condition.  This is a third declared criterion.  It must not be silently identified
with literal set intersection.  The exact intersection proposition, the exact set
distance, and the tolerance-level decision remain distinct even when they are reported
together.

\section{The model hierarchy}

The ICMSEP problem is indexed by a hierarchy of candidate classes rather than by one
unconstrained matrix space.  A lower-order class may impose a fixed orbital basis,
short hopping range, or restricted symmetry parameterization; a larger class may add
neighbor shells or selected symmetry-allowed corrections.  The scientific decision is
the first tested class for which the declared compatibility predicate passes.

Nested model classes do not automatically make every computed loss monotone.  Enlarging
$\mathfrak M_j$ preserves existential feasibility only when the class inclusion, the
reference data, the losses, the tolerances, and the alignment-family embedding are
conformant.  Refitting, changing an optimizer basin, changing support, or changing the
alignment family can break naive monotonicity of reported minimizers.  The finite
first-passage logic developed in Chapter~\ref{ch:reduction-decision} therefore examines
every earlier declared class directly rather than assuming monotone numerical behavior.

The phrase ``smallest adequate model'' is consequently shorthand for a precise finite
statement: the selected class passes the joint predicate, and every earlier class in the
declared inventory fails it.  It is not a claim about all conceivable Hamiltonian
families.

\section{The alignment problem is a component of operator admissibility}

The operator loss cannot be evaluated until the reference and candidate coordinates
have been identified.  For momentum-indexed retained operators, the aligned difference
used in this companion is
\[
  \Delta_{C,\theta}(k)
  =H_{\mathrm{ref}}(k)
   -C(k)^{-1}H_j(\theta;k)C(k).
\]
A weighted Frobenius operator loss may then be formed from momentum-space differences or
from corresponding real-space blocks.  The exact choice of weights, blocks, and
normalization belongs to the definition of $\mathcal L_{H,j}$.

The admissible alignment family $\mathcal U_j$ is part of that definition.  Identity
alignment, one global unitary, arbitrary pointwise unitaries, and symmetry-constrained
or finite-depth families produce different operator-admissible sets.  Enlarging the
family can reduce the optimized operator loss, but it can also admit transformations
that alter Fourier support or violate an application's structural requirements.

Thus the family-relative question
\[
  \exists C\in\mathcal U_j,
  \qquad
  \mathcal L_{H,j}(C,\theta)\leq\tau_H
\]
is a fiber of the full compatibility problem at one candidate $\theta$.  It is not a
replacement for the joint quantification over candidate Hamiltonians.  Chapters
\ref{ch:finite-operator-foundations} and \ref{ch:bloch-alignment-locality} develop this
operator branch because it is necessary for defining $\mathfrak A_{H,j}$ correctly.

\section{Fourier locality is a structural diagnostic}

For a momentum-dependent retained Hamiltonian, finite Fourier transformation produces
real-space blocks.  Hopping range, blockwise residuals, and truncation error are
therefore properties of the operator representation.  A unitary common to all momentum
points passes through the Fourier sum, while a momentum-dependent unitary can
redistribute Fourier weight among modes.  Consequently, a pointwise alignment that
improves the operator loss may change locality diagnostics.

The Fourier results in this companion establish three bounded facts.  Fixed-unitary
conjugation commutes with the finite Fourier operations used here.  An exact two-point
rank-two construction disproves universal commutation between pointwise conjugation and
support truncation.  Exact finite-tail identities and inequalities then control the
deliberate-truncation channel for a declared support.

These results clarify how $\mathcal L_{H,j}$ and a locality or support requirement may
interact.  They do not define the spectral admissible set, select a material model, or
replace the full spectral--operator intersection problem.

\section{Four interacting mathematical layers}

The proof architecture can now be stated in its proper order.

\subsection{Candidate and criterion layer}

A declared model class, reference problem, spectral loss, operator loss, tolerances, and
alignment family determine the two admissible sets.  Their intersection is the primary
compatibility proposition.  Their metric separation quantifies an obstruction under
additional geometric hypotheses.

\subsection{Representation and alignment layer}

Finite frames represent retained subspaces.  Projectors are invariant under internal
right-unitary changes, compressed operators transform covariantly, and reconstructed
ambient operators are invariant.  Between two models, an admissible alignment pulls the
candidate into reference coordinates before an operator residual is formed.

\subsection{Fourier and locality layer}

Finite Fourier transformation converts the aligned operator field into real-space
blocks.  Fixed and pointwise coordinate changes have different support consequences.
For a reference field, an interpolant, and its deliberately truncated approximation,
the interpolation, truncation, and total residuals form an exact additive chain at the
vector level.  Their norms satisfy triangle bounds rather than a scalar additive
identity.  Only the deliberate-truncation leg is identified with the checked omitted
Fourier tail.

\subsection{Geometry, certificate, and decision layer}

Exact admissible-set geometry proves intersection or separation for the families covered
by its hypotheses.  Certificate checkers connect finite data to exact propositions.
Three-way interval decisions preserve unresolved threshold overlap, and finite
first-passage logic selects the earliest class in one declared inventory.

\section{Organization}

Chapter~\ref{ch:finite-operator-foundations} proves the representation-neutral frame,
projector, compression, and reconstruction identities needed by the aligned operator
criterion.  Chapter~\ref{ch:bloch-alignment-locality} adds momentum dependence, aligned
operator differences, alignment families, fixed and pointwise gauge actions, the
three-channel residual calculus, the rank-two noncommutation witness, and exact
Fourier-tail control.
Chapter~\ref{ch:reduction-decision} returns to the full spectral--operator compatibility
problem and develops admissible-set geometry, separation, certificate semantics,
interval dispositions, and selection over a finite model inventory.
Chapter~\ref{ch:mechanization-status} states what has and has not yet been assembled into
an end-to-end formal adequacy theorem.

The appendices provide declaration-level traceability, immutable provenance, audit
boundaries, notation bridges, and a source- and status-qualified crosswalk of the
error-channel development across the seven-paper program.  They separate repository
mechanics from the mathematical narrative without changing the authority of the
underlying Lean statements.

\section{Difference convention and scope}

Throughout the companion, the aligned difference is reference minus pulled-back
candidate:
\[
  \Delta_{\mathrm{ref},\mathrm{cand}}(k)
  =H_{\mathrm{ref}}(k)-C(k)^{-1}H_{\mathrm{cand}}(k)C(k).
\]
Reversing the operands negates the difference.  Norm-based losses are insensitive to
this sign, but signed perturbations, defect potentials, and response coefficients are
not.

All results concern exact finite objects under explicit hypotheses.  They do not by
themselves establish that a runtime frame is exactly orthonormal, that floating-point
intervals contain exact quantities, that a sampled spectral gap persists between sample
points, or that a candidate is physically adequate for silicon.  In particular, the
companion supplies formal components of the ICMSEP decision architecture; it does not
claim completion of the proposed first-principles calculation, a retained certificate,
or a material-validation result.

\chapter{Finite-Dimensional Operator Foundations}
\label{ch:finite-operator-foundations}

\section{The finite retained-space problem}

Finite-dimensional reduction replaces a large operator problem by its action on a
smaller declared subspace. Let $\mathcal H_{\mathrm{amb}}\cong\mathbb C^D$ be an ambient
Hilbert space carrying an operator $H$, and let
$\mathcal H_{\mathrm{ret}}\cong\mathbb C^M$, with $M\leq D$, be a retained coordinate
space. A rectangular matrix $V\in\mathbb C^{D\times M}$ maps retained-coordinate vectors
into the ambient space. Its columns may be selected eigenvectors, projected trial
states, reduced-basis vectors, or any other declared finite frame.

Two mathematical objects must be kept distinct. The range of $V$ is the retained
subspace inside the ambient Hilbert space. The ordered columns of $V$ are one coordinate
frame for that subspace. When the columns are orthonormal, the subspace is represented
in ambient coordinates by
\[
  P_V=VV^\dagger,
\]
whereas the matrix of $H$ in retained coordinates is
\[
  C_V(H)=V^\dagger HV.
\]
The first object answers which ambient vectors are retained. The second answers how the
ambient operator acts after it has been expressed in a selected retained basis.

The distinction matters because a retained subspace has no unique ordered basis. If
$G\in U(M)$, then $VG$ spans the same subspace as $V$. A sound reduction formalism must
therefore show that the ambient projector is unchanged, identify the covariance law for
the compressed matrix, and verify that reconstruction in the ambient space does not
depend on the retained-coordinate basis. Otherwise a change of coordinates can be
mistaken for a change of model.

The ten theorems in this chapter establish those finite identities without invoking
periodicity, crystal momentum, Fourier transformation, or spatial locality. They form a
representation-neutral foundation that can be reused for Bloch frames, trial subspaces,
finite quantum systems, and other matrix reductions. The later Bloch chapter adds the
extra structure carried by a momentum-indexed family of frames.

\section{Definitions and notation}

Let the ambient and retained Hilbert spaces be
\[
  \mathcal H_{\mathrm{amb}}\cong\mathbb C^D,
  \qquad
  \mathcal H_{\mathrm{ret}}\cong\mathbb C^M,
  \qquad M\leq D.
\]
The Lean declarations use arbitrary finite index types rather than natural-number
indices; $D$ and $M$ denote their cardinalities in the standard matrix notation used
here.  The modeling notation assumes $M\leq D$, but the Lean definitions and theorem
statements impose no separate cardinality inequality; exact frame orthonormality appears
as an explicit equation wherever it is needed.  The symbol $\dagger$ denotes conjugate
transpose, and $I_D$ and $I_M$ denote the identity matrices on the ambient and retained
spaces.

In Dirac notation, choose the coordinate basis
$\{|a\rangle_{\mathrm{ret}}\}_{a=1}^M$ of the retained space and let
$\{|v_a\rangle\}_{a=1}^M$ be the retained vectors in the ambient space. The frame matrix
is the linear embedding
\[
  V=\sum_{a=1}^M |v_a\rangle\,{}_{\mathrm{ret}}\!\langle a|,
  \qquad
  V^\dagger=\sum_{a=1}^M |a\rangle_{\mathrm{ret}}\langle v_a|.
\]
Thus $V$ is not one state when $M>1$; it is the ordered collection of retained kets,
viewed as a map from retained coordinates into the ambient Hilbert space. The identities
\[
  V^\dagger V=I_M,
  \qquad
  VV^\dagger=\sum_{a=1}^M |v_a\rangle\langle v_a|
\]
express, respectively, orthonormality of the retained kets and the associated ambient
projector. The compressed operator has matrix elements
\[
  \langle a|C_V(H)|b\rangle
  =\langle v_a|H|v_b\rangle.
\]
A right-unitary change $V\mapsto VG$ replaces the retained kets by
\[
  |v'_a\rangle=\sum_{b=1}^M |v_b\rangle G_{ba}.
\]
For $M=1$, the embedding reduces to a single normalized ket, $V=|v\rangle$, with
$P_V=|v\rangle\langle v|$ and $C_V(H)=\langle v|H|v\rangle$.

The variables used throughout this chapter are:

\begin{itemize}
  \item $V\in\mathbb C^{D\times M}$: a finite retained frame. Its columns are the
        retained vectors expressed in the ambient basis.
  \item $H\in\mathbb C^{D\times D}$: an ambient operator, typically interpreted as a
        finite Hamiltonian when $H^\dagger=H$.
  \item $G\in\mathbb C^{M\times M}$: a retained-coordinate unitary satisfying
        $G^\dagger G=GG^\dagger=I_M$.
  \item $V_G=VG$: the same frame after a right-unitary change of retained coordinates.
  \item $P_V=VV^\dagger\in\mathbb C^{D\times D}$: the ambient frame operator. When the
        columns of $V$ are orthonormal, $P_V$ is the orthogonal projector onto their
        span.
  \item $C_V(H)=V^\dagger HV\in\mathbb C^{M\times M}$: the compression of $H$ to
        retained coordinates.
  \item $R_V(H)=V C_V(H)V^\dagger\in\mathbb C^{D\times D}$: reconstruction of the
        compressed operator in the ambient coordinates.
\end{itemize}

\begin{definition}[Orthonormal finite frame]
The frame $V$ is orthonormal when
\[
  V^\dagger V=I_M.
\]
This is an exact hypothesis. It is not inferred from a floating-point tolerance.
\end{definition}

\begin{definition}[Hermitian ambient operator]
The operator $H$ is Hermitian when $H^\dagger=H$.
\end{definition}

\begin{remark}[Formal correspondence]
In \modulepath{ProjectKoios/BandGauge/Alignment/FiniteFrame.lean}, the definitions above
are named \decl{IsOrthonormalFrame}, \decl{frameProjector},
\decl{rightUnitaryFrame}, \decl{compressOperator}, and
\decl{reconstructCompressedOperator}. The prose proofs below explain the checked Lean
theorems but do not replace their proof terms.
\end{remark}

\section{Relation to the research manuscript's finite-dimensional foundation}

The working \texttt{ksdft2effmass} research manuscript uses the same representation-neutral
objects in its finite-dimensional-foundations and projection-and-retained-spaces
chapters. Table~\ref{tab:finite-manuscript-notation-bridge} records the notation bridge.

\begin{table}[ht]
\centering
\small
\begin{tabularx}{\textwidth}{@{}l X X@{}}
\toprule
This paper & Research manuscript & Common mathematical object \\
\midrule
$V$
  & $V$
  & $D\times M$ retained-frame embedding \\
$|v_a\rangle$
  & $|\phi_a\rangle$
  & column $a$ of the retained frame \\
$P_V=VV^\dagger$
  & $P$ or $\widehat P_{\mathcal S}$
  & ambient projector onto the retained subspace \\
$C_V(H)=V^\dagger HV$
  & $H_{\mathcal S}$
  & operator matrix in retained coordinates \\
$R_V(H)=VC_V(H)V^\dagger$
  & $PHP$ embedded in the parent space
  & reconstructed retained operator in ambient coordinates \\
$V_G=VG$
  & $V'=VG$
  & right-unitary change of retained frame \\
\bottomrule
\end{tabularx}
\caption{Representation-neutral notation bridge to the working research manuscript.}
\label{tab:finite-manuscript-notation-bridge}
\end{table}

The theorem correspondence is likewise direct. Declarations
\decl{BG-ALIGNMENT-THM-001}, \decl{BG-ALIGNMENT-THM-002},
\decl{BG-ALIGNMENT-THM-009}, and \decl{BG-ALIGNMENT-THM-010} justify the projector
properties used in the projection discussion. Declarations
\decl{BG-ALIGNMENT-THM-003} and \decl{BG-ALIGNMENT-THM-004} justify right-unitary frame
changes and projector invariance; \decl{BG-ALIGNMENT-THM-004} realizes the formula
presented prospectively as \texttt{PRF-05.01}. Declaration
\decl{BG-ALIGNMENT-THM-005} realizes compression covariance, presented prospectively as
\texttt{PRF-05.02}. The remaining results establish Hermiticity of compression,
projector-sandwich reconstruction, and reconstruction invariance.

This bridge is editorial, not authoritative. It records notation observed in the clean
\texttt{ksdft2effmass} worktree on branch
\decl{work/monograph-representation-carriers-and-derivation-graphs} at commit
\decl{e3d8a51b30e2072aa4bb116eb46389e342fc00fc}, principally in
\modulepath{chapters/finite-dimensional-foundations.tex} and
\modulepath{chapters/projection-and-retained-spaces.tex}. Those chapters supply the
source notation for this bridge. The Lean declarations identified here fix the exact
statements and proof status.

\section{Finite-frame theorem sequence}

\begin{theorem}[\texttt{BG-ALIGNMENT-THM-001}: Hermiticity of the frame operator]
\label{thm:bg-alignment-001}
For every finite complex frame $V\in\mathbb C^{D\times M}$,
\[
  P_V^\dagger=P_V.
\]
No orthonormality hypothesis is required.
\end{theorem}

\begin{proof}
By conjugate-transpose reversal and involution,
\[
  P_V^\dagger=(VV^\dagger)^\dagger
  =(V^\dagger)^\dagger V^\dagger
  =VV^\dagger=P_V.
\]
\end{proof}

Hermiticity is the first structural requirement for interpreting $P_V$ as an observable-
compatible subspace operator. Importantly, it holds before normalization is imposed.
This separates two logically different facts: every frame produces a Hermitian Gram-
type operator in the ambient space, whereas only an orthonormal frame makes that
operator an orthogonal projector. The distinction prevents approximate normalization
from being hidden inside a symmetry claim.

\begin{theorem}[\texttt{BG-ALIGNMENT-THM-002}: Idempotence of an orthonormal-frame projector]
\label{thm:bg-alignment-002}
If $V^\dagger V=I_M$, then
\[
  P_V^2=P_V.
\]
Together with Theorem~\ref{thm:bg-alignment-001}, this makes $P_V$ the orthogonal
projector onto the retained column space.
\end{theorem}

\begin{proof}
Associativity exposes the retained-space Gram matrix:
\[
  P_V^2=(VV^\dagger)(VV^\dagger)
  =V(V^\dagger V)V^\dagger
  =VI_MV^\dagger
  =VV^\dagger=P_V.
\]
\end{proof}

Together, Hermiticity and idempotence identify $P_V$ as the coordinate-free carrier of
the retained subspace. This is the object that should be compared when the scientific
question concerns which states are retained rather than how a basis for those states is
written. The theorem is also the point at which exact orthonormality becomes essential;
without it, $VV^\dagger$ need not be a projector.

\begin{theorem}[\texttt{BG-ALIGNMENT-THM-003}: Orthonormality after a retained unitary]
\label{thm:bg-alignment-003}
If $V^\dagger V=I_M$ and $G$ is unitary, then $V_G=VG$ is also an orthonormal frame:
\[
  V_G^\dagger V_G=I_M.
\]
\end{theorem}

\begin{proof}
Using conjugate-transpose reversal, the frame hypothesis, and $G^\dagger G=I_M$,
\[
  V_G^\dagger V_G
  =(VG)^\dagger(VG)
  =G^\dagger(V^\dagger V)G
  =G^\dagger I_MG
  =I_M.
\]
\end{proof}

Composite retained states admit internal unitary mixing, and numerical eigensolvers may
return different orthonormal bases for the same retained subspace. This theorem shows
that such a mixing remains inside the class of admissible orthonormal frames. Smoothness,
optimization, and physical selection enter only when a particular gauge is chosen.

\begin{theorem}[\texttt{BG-ALIGNMENT-THM-004}: Projector invariance under a retained unitary]
\label{thm:bg-alignment-004}
For every frame $V$ and every retained-space unitary $G$,
\[
  P_{V_G}=P_V.
\]
The conclusion does not require $V^\dagger V=I_M$.
\end{theorem}

\begin{proof}
The unitary and its adjoint cancel in the retained coordinates:
\[
  P_{V_G}=(VG)(VG)^\dagger
  =VGG^\dagger V^\dagger
  =VI_MV^\dagger
  =P_V.
\]
\end{proof}

The projector is unchanged even though the retained basis has changed. Consequently,
$V$ and $VG$ describe the same ambient subspace, not two competing physical models.
This invariance is the basic reason that a subspace comparison should not penalize a
pure retained-coordinate rotation. It also explains why projector agreement is weaker
than entrywise agreement of compressed operators.

\begin{theorem}[\texttt{BG-ALIGNMENT-THM-005}: Covariance of operator compression]
\label{thm:bg-alignment-005}
For every ambient operator $H$, frame $V$, and retained-space unitary $G$,
\[
  C_{V_G}(H)=G^\dagger C_V(H)G.
\]
Neither Hermiticity of $H$ nor orthonormality of $V$ is required.
\end{theorem}

\begin{proof}
Expanding the changed frame and reversing the conjugate transpose of its product gives
\[
  C_{V_G}(H)
  =(VG)^\dagger H(VG)
  =G^\dagger V^\dagger HVG
  =G^\dagger C_V(H)G.
\]
\end{proof}

A compressed operator is not invariant as a matrix: it transforms by unitary
similarity. Thus its individual entries are coordinate-dependent, while similarity-
invariant information such as its characteristic polynomial is not. This covariance is
why operator comparisons must first place both reduced matrices in a common retained
coordinate system. Selection of the aligning unitary $G$ is a separate problem.

\begin{theorem}[\texttt{BG-ALIGNMENT-THM-006}: Hermiticity of compression]
\label{thm:bg-alignment-006}
If $H^\dagger=H$, then for every frame $V$,
\[
  C_V(H)^\dagger=C_V(H).
\]
Frame orthonormality is not required for this algebraic closure.
\end{theorem}

\begin{proof}
Conjugate-transpose reversal and the Hermiticity of $H$ yield
\[
  C_V(H)^\dagger
  =(V^\dagger HV)^\dagger
  =V^\dagger H^\dagger V
  =V^\dagger HV
  =C_V(H).
\]
\end{proof}

When $H$ is a finite Hamiltonian, this result ensures that its compressed matrix remains
Hermitian and can therefore be treated as a valid finite Hermitian eigenproblem. The
closure is algebraic. Spectral approximation and subspace adequacy require additional
hypotheses.

\begin{theorem}[\texttt{BG-ALIGNMENT-THM-007}: Reconstruction is a projector sandwich]
\label{thm:bg-alignment-007}
For every frame $V$ and ambient operator $H$,
\[
  R_V(H)=P_VHP_V.
\]
Here $P_V$ need not be idempotent unless the frame is orthonormal.
\end{theorem}

\begin{proof}
This is exact reassociation of the matrix products:
\[
  R_V(H)
  =V(V^\dagger HV)V^\dagger
  =(VV^\dagger)H(VV^\dagger)
  =P_VHP_V.
\]
\end{proof}

The identity $R_V(H)=P_VHP_V$ exposes the ambient operator actually retained by the
finite reduction. It shows that compression followed by reconstruction is a two-sided
restriction to the chosen frame operator, rather than an unexplained smaller matrix.
For an orthonormal frame this is the familiar projected operator $P_VHP_V$; the theorem
itself carefully avoids assuming that $P_V$ is already idempotent.

\begin{theorem}[\texttt{BG-ALIGNMENT-THM-008}: Reconstruction invariance]
\label{thm:bg-alignment-008}
For every frame $V$, retained-space unitary $G$, and ambient operator $H$,
\[
  R_{V_G}(H)=R_V(H).
\]
\end{theorem}

\begin{proof}
Apply Theorem~\ref{thm:bg-alignment-007} to the changed frame and then use
Theorem~\ref{thm:bg-alignment-004}:
\[
  R_{V_G}(H)=P_{V_G}HP_{V_G}=P_VHP_V=R_V(H).
\]
Thus the retained matrix changes by unitary similarity, but its reconstructed ambient
operator does not change.
\end{proof}

Although the compressed matrix changes under $V\mapsto VG$, the ambient reconstructed
operator does not. This is the precise sense in which retained-coordinate gauge changes
alter representation but not the reconstructed reduced operator. It provides a
basis-independent comparison object and guards against treating a similarity transform
as a physical discrepancy.

\begin{theorem}[\texttt{BG-ALIGNMENT-THM-009}: The projector fixes its frame]
\label{thm:bg-alignment-009}
If $V^\dagger V=I_M$, then
\[
  P_VV=V.
\]
\end{theorem}

\begin{proof}
By associativity and the exact frame hypothesis,
\[
  P_VV=(VV^\dagger)V=V(V^\dagger V)=VI_M=V.
\]
\end{proof}

The projector acts as the identity on every retained column. This confirms that the
ambient projection leaves the represented retained states unchanged, a basic
consistency condition for interpreting those columns as the range of $P_V$. On the orthogonal
complement, the projector acts as zero rather than as the ambient identity.

\begin{theorem}[\texttt{BG-ALIGNMENT-THM-010}: Compression of the retained projector]
\label{thm:bg-alignment-010}
If $V^\dagger V=I_M$, then the compression of $P_V$ is the retained identity:
\[
  C_V(P_V)=I_M.
\]
\end{theorem}

\begin{proof}
Expanding $P_V$ and reassociating gives
\[
  C_V(P_V)
  =V^\dagger(VV^\dagger)V
  =(V^\dagger V)(V^\dagger V)
  =I_MI_M=I_M.
\]
\end{proof}

In retained coordinates, the projector becomes exactly $I_M$. This supplies a useful
normalization and regression identity: a correct compression of the retained projector
cannot introduce a nontrivial retained-space operator. It closes the relationship
between the ambient projector and its retained coordinate representation.

\begin{remark}[Scope of the theorem sequence]
These ten theorems concern exact finite matrices and equal-rank right-unitary coordinate
changes. Runtime frame construction, unequal-rank comparison, spectral adequacy, and
momentum-space regularity enter later layers with additional hypotheses.
\end{remark}

\chapter{Bloch Frames, Alignment Families, and Fourier Locality}
\label{ch:bloch-alignment-locality}

\section{The comparison problem in finite band reduction}

\subsection{From a parent Hamiltonian to a retained model}

Many electronic-structure and reduced-order calculations begin with an operator that is
too large, too detailed, or expressed in the wrong coordinates for the intended
analysis. On a finite momentum mesh, let $\mathcal K$ denote the supplied set of crystal
momenta and let $N_k=|\mathcal K|$. At each $k\in\mathcal K$, a parent calculation
provides a finite Hermitian matrix
\[
  H_{\mathrm{parent}}(k)\in\mathbb C^{D\times D}.
\]
The dimension $D$ may count plane waves, localized orbitals, finite-difference degrees
of freedom, or any other finite ambient basis; the same algebra applies to each case.

Suppose the scientific question concerns only $M$ states, with $M\leq D$, such as an
isolated band, a composite group of bands near the Fermi level, or a trial subspace used
for downfolding. These states are assembled as the columns of a frame
\[
  V(k)\in\mathbb C^{D\times M}.
\]
When $V(k)^\dagger V(k)=I_M$, the columns are orthonormal and determine the retained
subspace
\[
  \mathcal S(k)=\operatorname{Ran}V(k)
\]
and its ambient orthogonal projector
\[
  P(k)=V(k)V(k)^\dagger.
\]
The parent operator represented in retained coordinates is
\[
  h(k)=V(k)^\dagger H_{\mathrm{parent}}(k)V(k)
  \in\mathbb C^{M\times M}.
\]
This matrix is the finite projected Hamiltonian. Spectral isolation and accuracy for a
target observable are additional properties of the selected subspace.

In applications, the retained frame $V(k)$ may be obtained by diagonalizing the parent
Hamiltonian, projecting trial orbitals, disentangling a composite band manifold,
constructing Wannier functions, or fitting a reduced model. The finite-frame theorem
sequence in Chapter~\ref{ch:finite-operator-foundations} takes the resulting matrix
$V(k)$ and explicit hypotheses such as $V(k)^\dagger V(k)=I_M$ as its inputs.

\subsection{Why two reduced matrices cannot be compared entry by entry}

The columns of $V(k)$ are coordinates, not uniquely labeled physical objects. If
$G(k)\in U(M)$ is unitary, then
\[
  V(k)\longmapsto V(k)G(k)
\]
changes the retained basis but not the subspace $\mathcal S(k)$. The projector remains
$P(k)$, while the compressed Hamiltonian becomes
\[
  h(k)\longmapsto G(k)^\dagger h(k)G(k).
\]
Thus the subspace is invariant, but the retained matrix is only covariant. Its entries
and eigenvectors depend on the chosen retained coordinates even though its spectrum is
unchanged.

This distinction becomes unavoidable when comparing a reference model with a candidate
model. Let
\[
  h_{\mathrm{ref}}(k),\qquad h_{\mathrm{cand}}(k)
  \in\mathbb C^{M\times M}
\]
denote their retained operators. A direct difference
$h_{\mathrm{ref}}(k)-h_{\mathrm{cand}}(k)$ is meaningful only if the two matrices are
already expressed in the same retained coordinates. More generally, an identification
$U(k)$ from candidate coordinates to reference coordinates must be supplied, and the
comparison uses the pulled-back candidate operator
\[
  U(k)^{-1}h_{\mathrm{cand}}(k)U(k).
\]
The sign convention used in this repository is
\[
  \Delta_U(k)
  =h_{\mathrm{ref}}(k)
   -U(k)^{-1}h_{\mathrm{cand}}(k)U(k).
\]
For a unitary identification, $U(k)^{-1}=U(k)^\dagger$. The norm of $\Delta_U(k)$
measures an operator discrepancy relative to that supplied identification. Its
interpretation therefore includes the construction and admissibility of $U$.

Band energies are insufficient for this comparison. Pointwise unitary conjugation preserves
the eigenvalues of $h(k)$, so two fields can have identical band energies while their
matrix-valued eigenvectors, Fourier coefficients, and real-space representations differ.
An adequacy study that cares about operator structure must therefore specify more than a
spectral tolerance.

\subsection{Why momentum dependence affects locality}

A real-space reduced model is obtained by Fourier transforming the momentum-dependent
operator field. In conventional finite-mesh notation, its block at translation $R$ is
\[
  T_R[h]
  =\frac{1}{N_k}\sum_{k\in\mathcal K}
    e^{-\mathrm{i}k\cdot R}h(k).
\]
Here $R$ labels a real-space lattice translation; in the exact Lean development the
phase factors are finite-group characters. Short-range or sparse effective models keep
only a declared set of $R$ values. Consequently, Fourier support and the distribution of
Frobenius power among real-space blocks are part of the model-reduction question, not
merely presentation choices.

If one constant unitary $G$ is used at every $k$, it factors through the finite Fourier
sum:
\[
  T_R[G^\dagger hG]=G^\dagger T_R[h]G.
\]
The real-space labels are not mixed. Exact support truncation therefore commutes with
this common change of retained coordinates, and unitary-invariant block norms are
preserved.

The situation changes when $G$ depends on $k$:
\[
  T_R[G(\cdot)^\dagger h(\cdot)G(\cdot)]
  =\frac{1}{N_k}\sum_{k\in\mathcal K}
    e^{-\mathrm{i}k\cdot R}G(k)^\dagger h(k)G(k).
\]
The matrices $G(k)$ can no longer be factored outside the sum. A pointwise alignment may
therefore improve an operator comparison at each momentum while changing how the same
field is distributed over real-space blocks. This is the tension at the center of the
paper: pointwise spectral or operator agreement and real-space locality are not the same
criterion. Section~\ref{sec:gauge-truncation-counterexample} gives the checked exact
two-point witness: a varying rank-two gauge creates a nonzero odd mode and does not
commute with zero-mode support truncation.

\subsection{The admissible alignment family is part of the claim}

The phrase ``optimize over the gauge'' is incomplete until the allowed gauges are
specified. Several mathematically different questions arise:

\begin{itemize}
  \item Is no coordinate adjustment allowed, so that $U(k)=I_M$?
  \item May one global unitary be used at every momentum?
  \item May an arbitrary unitary be chosen independently at every momentum?
  \item Must $U(k)$ satisfy continuity, endpoint sewing, symmetry, orientation, or a
        finite-depth structural ansatz?
\end{itemize}

These choices define different feasible sets. Enlarging the alignment family can only
make an existential adequacy test easier to satisfy, while the reverse implication
fails in general. A pointwise witness belongs to the pointwise family; transfer to a
global, smooth, or structured family requires an explicit membership proof. The strict
family-separation theorems later in this chapter make this distinction exact.

For that reason, the operative adequacy question is joint and family-relative: does one
candidate, one support choice, and one alignment from the declared family satisfy all
operator, spectral, and locality thresholds? Such a witness uses the same alignment for
every objective. Envelopes, plots, and optimizer traces summarize searches rather than
supply the existential witness or an infeasibility certificate.

\subsection{Role of the finite-frame theorems}

The ten finite-frame results proved in
Chapter~\ref{ch:finite-operator-foundations} supply the algebraic foundation for this
comparison. They show exactly which objects are invariant when only retained coordinates
change, which objects transform covariantly, and how a compressed operator can be
reconstructed as an ambient projector sandwich. This prevents three common category
errors:

\begin{enumerate}
  \item treating a basis change as a change of retained subspace;
  \item comparing compressed matrices before identifying their retained coordinates;
        and
  \item inferring physical or spectral adequacy from algebraic covariance alone.
\end{enumerate}

The present statements are finite and exact. Construction of $V(k)$ and $U(k)$,
continuity, localization, and material interpretation enter through additional
mathematical or evidential layers.

\section{Relation to the research manuscript's Bloch and operator-comparison chapters}

The working composite-band chapter denotes its momentum-dependent frame by
\[
  F(\kappa)=\begin{bmatrix}u_1(\kappa)&\cdots&u_M(\kappa)\end{bmatrix}.
\]
At corresponding momentum coordinates, the notation of this paper is
\[
  V(k)=F(\kappa)
  =\sum_{a=1}^M |u_a(\kappa)\rangle\,{}_{\mathrm{ret}}\!\langle a|.
\]
Thus $V(k)$ and $F(\kappa)$ are two names for the map that inserts retained-coordinate
vectors into an ambient Bloch fiber.

\begin{table}[ht]
\centering
\small
\begin{tabularx}{\textwidth}{@{}l X X@{}}
\toprule
This paper & Research manuscript & Common mathematical object \\
\midrule
$V(k)$
  & $F(\kappa)$
  & momentum-dependent retained-frame embedding \\
$|v_a(k)\rangle$
  & $u_a(\kappa)$
  & column $a$ of the composite Bloch frame \\
$P_V(k)$
  & $P(\kappa)=F(\kappa)F(\kappa)^\dagger$
  & ambient projector field \\
$C_{V(k)}(H(k))$
  & $H_F(\kappa)$
  & retained operator field \\
$T_R[h]$
  & $H_U(R)$ or $\langle w_{m0}|\widehat H|w_{nR}\rangle$
  & real-space Fourier block of the retained operator \\
\bottomrule
\end{tabularx}
\caption{Bloch and Fourier notation bridge to the working research manuscript.}
\label{tab:bloch-manuscript-notation-bridge}
\end{table}

This paper uses $G(k)$ for an internal unitary change of a single retained frame and
reserves $U(k)$ for an identification between reference and candidate retained spaces.
The working composite-band chapter writes its internal gauge as $U(\kappa)$. These
uses must not be merged silently: an internal gauge and a two-model identification have
different domains, provenance, and adequacy obligations.

The aligned-difference sign also differs. When the research manuscript's pristine
operator is identified with $h_{\mathrm{ref}}$, its doped operator with
$h_{\mathrm{cand}}$, and its map $U$ runs from reference to candidate coordinates, it
writes
\[
  \Delta H_{d\to b}=U^\dagger h_{\mathrm{cand}}U-h_{\mathrm{ref}}.
\]
The Formal convention used here is
\[
  \Delta_U=h_{\mathrm{ref}}-U^{-1}h_{\mathrm{cand}}U
  =-\Delta H_{d\to b}
\]
for unitary $U$. Norms are insensitive to this sign, but a signed perturbation or
impurity interpretation is not. Any later editorial export must state the operand order.

This bridge records notation observed at the same clean research-manuscript commit,
principally in \modulepath{chapters/composite-band-gauge-alignment.tex} and
\modulepath{chapters/operator-comparison.tex}. Formal source fixes theorem status.

\section{Aligned operator differences}

The Bloch layer represents a pointwise unitary alignment by an isometric linear
equivalence $U(k)$ between reference and candidate fibers. The candidate operator is
pulled back into the reference fiber before subtraction. Declarations
\decl{BG-BLOCH-THM-012--020} prove pointwise evaluation, closure of shared sewing under
subtraction, sign reversal, recovery, an exact-zero criterion, and sewing closure of the
aligned difference.

The sewing theorem assumes sewing laws for both operator fields and an alignment that
intertwines them. Under these hypotheses, the aligned difference inherits sewing.
Self-adjointness and regularity require their own hypotheses.

\section{Alignment families are decision spaces}

Declarations \decl{BG-ALIGNMENT-DEF-007--016} distinguish:

\begin{itemize}
  \item one declared singleton alignment;
  \item identity-only alignment;
  \item one global unitary reused at every point; and
  \item the complete family of arbitrary pointwise unitary alignments.
\end{itemize}

Feasibility means existence of an accepted alignment inside one specified family.
Declarations \decl{BG-ALIGNMENT-THM-011--019} prove the forward inclusions
\[
  \mathcal U_{\mathrm{id}}
  \subseteq
  \mathcal U_{\mathrm{global}}
  \subseteq
  \mathcal U_{\mathrm{pointwise}}
\]
and transfer feasibility only in that direction.

Declarations \decl{BG-ALIGNMENT-THM-020--029} provide strictness witnesses when the
relevant point and unitary inequalities are supplied. They also construct acceptance
predicates feasible in a larger family but infeasible in a proper subfamily. Therefore a
pointwise-alignment witness need not belong to the global-unitary or identity family.
The result is a logical separation of decision spaces.

\section{Fixed and pointwise gauge actions}

A single unitary common to all finite coordinates preserves Frobenius inner products,
powers, and differences and commutes with finite Fourier coefficient extraction,
reconstruction, exact support truncation, and the corresponding residual. These results
are \decl{BG-FOURIER-THM-033--047}.

A pointwise unitary field has a different boundary. Declarations
\decl{BG-FOURIER-THM-048--054} prove an exact character-phase mode shift, the triviality
of scalar operator conjugation, coefficient separation, reduction of a constant field
to the fixed-unitary case, conditional coefficient change for a genuinely changed
field, and preservation of weighted pointwise Frobenius power.

The last result is invariance of weighted pointwise Frobenius power. Fourier
coefficients, support, decay, and locality require separate statements; a general
pointwise-gauge convolution theorem is outside the current family.

\section{An exact gauge--truncation counterexample}
\label{sec:gauge-truncation-counterexample}

Let the momentum group be $G=\mathbb Z/2\mathbb Z$, and define
\[
  P_0=\begin{bmatrix}1&0\\0&0\end{bmatrix},\qquad
  P_1=\begin{bmatrix}0&0\\0&1\end{bmatrix},\qquad
  S=\begin{bmatrix}0&1\\1&0\end{bmatrix}.
\]
The original rank-two field is constant, $H(k)=P_0$.  The pointwise unitary gauge is
$W(0)=I$ and $W(1)=S$.  Write $H^W(k)=W(k)^\dagger H(k)W(k)$, and let
$\mathcal T_0$ reconstruct a field after retaining only its trivial-character Fourier
coefficient.

\begin{theorem}[Two-point gauge--truncation noncommutation]
\label{thm:two-point-gauge-truncation}
For the exact field and gauge above, the original odd coefficient is zero, whereas
\[
  \widehat{H^W}_{\mathrm{odd}}
  =\frac12\begin{bmatrix}1&0\\0&-1\end{bmatrix}\neq 0.
\]
Moreover,
\[
  \mathcal T_0(H^W)\neq (\mathcal T_0 H)^W.
\]
\end{theorem}

\begin{proof}
Conjugation by $S$ exchanges $P_0$ and $P_1$.  Hence $H^W(0)=P_0$ and
$H^W(1)=P_1$.  The nontrivial character of $\mathbb Z/2\mathbb Z$ takes values $1$
and $-1$ at the two points, so the normalized odd coefficient is
$(P_0-P_1)/2$, while the same coefficient of the constant field vanishes.  Zero-mode
truncation fixes $H$.  If the two displayed routes were equal, evaluating their
zero-mode reconstruction at both momenta would force $P_0=P_1$, contradicting their
$(0,0)$ entries.  This is the checked proof sequence
\path{BG-FOURIER-THM-055}--\path{BG-FOURIER-THM-066}; the final declaration is
\path{ProjectKoios.BandGauge.Fourier.pointwiseSwapConjugation_zeroModeFourierTruncation_ne}.
\end{proof}

The theorem is a finite existence counterexample, not a universal statement about every
momentum-dependent gauge.  It establishes neither continuum locality nor decay,
smoothness, mesh convergence, runtime correctness, topology, or material adequacy.  In
rank one, scalar operator conjugation is trivial, so this rank-two witness must not be
reinterpreted as rank-one gauge-locality behavior.  The source module for the following
crosswalk is
\path{ProjectKoios/BandGauge/Fourier/GaugeTruncationCounterexample.lean}.

\begin{table}[ht]
\centering
\footnotesize
\begin{tabularx}{\textwidth}{@{}l >{\raggedright\arraybackslash}X@{}}
\toprule
ID & Exact Lean declaration and proof dossier \\
\midrule
\path{BG-FOURIER-THM-062}
  & \path{ProjectKoios.BandGauge.Fourier.pointwiseSwapConjugation_odd_coefficient}
    \newline \path{proofs/BandGauge/Fourier/BG-FOURIER-THM-062-gauged-odd-coefficient.tex} \\
\path{BG-FOURIER-THM-064}
  & \path{ProjectKoios.BandGauge.Fourier.pointwiseSwapConjugation_odd_coefficient_ne_zero}
    \newline \path{proofs/BandGauge/Fourier/BG-FOURIER-THM-064-gauged-odd-mode-nonzero.tex} \\
\path{BG-FOURIER-THM-066}
  & \path{ProjectKoios.BandGauge.Fourier.pointwiseSwapConjugation_zeroModeFourierTruncation_ne}
    \newline \path{proofs/BandGauge/Fourier/BG-FOURIER-THM-066-gauge-truncation-noncommutation.tex} \\
\bottomrule
\end{tabularx}
\caption{Key declarations for the exact counterexample.}
\label{tab:gauge-truncation-crosswalk}
\end{table}

\section{Exact finite truncation tails}
\label{sec:finite-truncation-tails}

Let $G$ be a finite additive group, let $\widehat G$ be its finite character group, and
let $c:\widehat G\to V$ be coefficients in a complex normed space.  For a declared
support $A\subseteq\widehat G$, write $c_A$ for the coefficients obtained by setting
all modes outside $A$ to zero, and define the deliberate-truncation residual
\[
  R_A(k)=\mathcal F^{-1}c(k)-\mathcal F^{-1}c_A(k).
\]
This is the interpolant-minus-truncated channel.  It is not the
reference-minus-interpolant channel and is not, by itself, the total residual.

\begin{theorem}[Finite Fourier-tail identities and bounds]
\label{thm:finite-fourier-tail-bounds}
Under the repository's finite Fourier normalization, the following hold.
\begin{enumerate}
  \item If $V$ is a complex Hilbert space, then
  \[
    \left(\frac1{|G|}\sum_{k\in G}\lVert R_A(k)\rVert^2\right)^{1/2}
    =\left(\sum_{\chi\notin A}\lVert c(\chi)\rVert^2\right)^{1/2}.
  \]
  \item For complex matrix coefficients, the same identity holds with the explicit
  entrywise Frobenius power
  $\sum_{i,j}|R_A(k)_{ij}|^2$ and
  $\sum_{i,j}|c(\chi)_{ij}|^2$.  The Lean root-mean-square theorems in items 1 and 2
  carry an explicit nonempty finite-domain premise.
  \item For any complex normed space and every $k\in G$,
  \[
    \lVert R_A(k)\rVert
    \leq \sum_{\chi\notin A}\lVert c(\chi)\rVert,
  \]
  and therefore
  \[
    \max_{k\in G}\lVert R_A(k)\rVert
    \leq \sum_{\chi\notin A}\lVert c(\chi)\rVert.
  \]
\end{enumerate}
\end{theorem}

\begin{proof}
The residual is definitionally the previously formalized finite Fourier truncation
residual.  Finite Parseval gives total residual power as $|G|$ times omitted coefficient
power; division by $\sqrt{|G|}$ yields the root-mean-square identities.  Pointwise, the
omitted-mode reconstruction formula and the triangle inequality bound the residual by
the sum of the norms of its summands.  Character phases disappear because each has norm
one.  The finite maximum statement is the coordinatewise supremum lift of this same
bound.  These steps are declarations \path{BG-FOURIER-THM-067}--
\path{BG-FOURIER-THM-075}.
\end{proof}

The matrix root-mean-square statement uses explicit entrywise Frobenius power, not the
ambient matrix norm.  Conversely, the maximum theorem uses the ambient norm on an
arbitrary complex normed space; no entrywise matrix-Frobenius maximum specialization is
claimed.  The maximum theorem requires no separate \path{Nonempty G} argument beyond
the additive-group structure, which already supplies $0:G$; this is not an empty-domain
claim.  None of the results binds floating-point coefficients, a runtime interpolant, a
mesh, physical hopping range, or a scientific threshold.  They also do not assert scalar
additivity of interpolation, truncation, and total errors.  The source module for both
following crosswalks is
\path{ProjectKoios/BandGauge/Fourier/ResidualChannels.lean}.

\begin{table}[ht]
\centering
\footnotesize
\begin{tabularx}{\textwidth}{@{}l >{\raggedright\arraybackslash}X@{}}
\toprule
ID & Exact Lean declaration and proof dossier \\
\midrule
\path{BG-FOURIER-THM-067}
  & \path{ProjectKoios.BandGauge.Fourier.truncationResidual_fourierReconstruction_eq}
    \newline \path{proofs/BandGauge/Fourier/BG-FOURIER-THM-067-truncation-channel.tex} \\
\path{BG-FOURIER-THM-068}
  & \path{ProjectKoios.BandGauge.Fourier.finiteFourierPower_truncationResidual_eq_omitted}
    \newline \path{proofs/BandGauge/Fourier/BG-FOURIER-THM-068-hilbert-tail.tex} \\
\path{BG-FOURIER-THM-069}
  & \path{ProjectKoios.BandGauge.Fourier.finiteMatrixFrobeniusPower_truncationResidual_eq_omitted}
    \newline \path{proofs/BandGauge/Fourier/BG-FOURIER-THM-069-matrix-tail.tex} \\
\path{BG-FOURIER-THM-070}
  & \path{ProjectKoios.BandGauge.Fourier.finiteRootMeanSquareNorm_truncationResidual_eq_sqrt_omitted}
    \newline \path{proofs/BandGauge/Fourier/BG-FOURIER-THM-070-hilbert-rms-tail.tex} \\
\path{BG-FOURIER-THM-071}
  & \path{ProjectKoios.BandGauge.Fourier.finiteRootMeanSquareNorm_truncationResidual_sq_eq_omitted}
    \newline \path{proofs/BandGauge/Fourier/BG-FOURIER-THM-071-hilbert-rms-squared-tail.tex} \\
\path{BG-FOURIER-THM-072}
  & \path{ProjectKoios.BandGauge.Fourier.finiteMatrixFrobeniusRootMeanSquare_truncationResidual_eq_sqrt_omitted}
    \newline \path{proofs/BandGauge/Fourier/BG-FOURIER-THM-072-matrix-frobenius-rms-tail.tex} \\
\path{BG-FOURIER-THM-073}
  & \path{ProjectKoios.BandGauge.Fourier.finiteMatrixFrobeniusRootMeanSquare_truncationResidual_sq_eq_omitted}
    \newline \path{proofs/BandGauge/Fourier/BG-FOURIER-THM-073-matrix-frobenius-rms-squared-tail.tex} \\
\bottomrule
\end{tabularx}
\caption{Exact power and root-mean-square declaration crosswalk.}
\label{tab:rms-tail-crosswalk}
\end{table}

\begin{table}[ht]
\centering
\footnotesize
\begin{tabularx}{\textwidth}{@{}l >{\raggedright\arraybackslash}X@{}}
\toprule
ID & Exact Lean declaration and proof dossier \\
\midrule
\path{BG-FOURIER-THM-074}
  & \path{ProjectKoios.BandGauge.Fourier.norm_truncationResidual_fourierReconstruction_le_omitted}
    \newline \path{proofs/BandGauge/Fourier/BG-FOURIER-THM-074-pointwise-maximum-tail.tex} \\
\path{BG-FOURIER-THM-075}
  & \path{ProjectKoios.BandGauge.Fourier.finiteMaximumNorm_truncationResidual_le_omitted}
    \newline \path{proofs/BandGauge/Fourier/BG-FOURIER-THM-075-finite-maximum-tail.tex} \\
\bottomrule
\end{tabularx}
\caption{Exact pointwise and finite-maximum declaration crosswalk.}
\label{tab:maximum-tail-crosswalk}
\end{table}
\FloatBarrier

\section{Finite route algebra}

The finite-lattice route layer defines coordinate pullback as $U^{\dagger}AU$.
Declarations \decl{SS-ROUTE-THM-001--010} prove compatibility with addition,
subtraction, composition, scalar-reference removal, Hermiticity, spin lift, folding, and
aligned defect extraction under their explicit isometry and identification premises.
They establish exact algebraic commutativity for the same declared maps. Nonlinear
fitting, separately selected gauges, and different model classes define different
routes and require separate analyses.

\chapter{Reduction, Certificates, and Finite Decisions}
\label{ch:reduction-decision}

\section{Admissible-set geometry}

For two subsets of one declared metric candidate space, the Formal geometry layer
separates a common witness from positive certified separation. Declarations
\decl{BG-THM-001--008} prove that a feasible pair bounds set distance above, a common
candidate forces zero distance, a positive lower bound forces disjointness, and
nonempty compact disjoint sets attain positive distance. Threshold-defined spectral and
operator sets are one application of these general facts.

A global separation conclusion requires a positive certificate on the declared
candidate space. The metric, losses, thresholds, and alignment family are part of that
claim; searches and plots are diagnostic inputs rather than separation proofs.

\section{Exact finite loss geometry}

The Hilbert and rank-two layers provide checked special cases rather than a universal
model-reduction theorem. \decl{BG-HILBERT-THM-001--015} cover exact ball geometry,
weighted Frobenius algebra, and finite spectral and operator residual bridges.
\decl{BG-THM-009--034} cover the canonical rank-two loss floor, ellipse sublevel sets,
intersection criteria, exact distances, and explicit rectangle-containment transfer.

The spectral perturbation declarations \decl{BG-FOUNDATION-THM-008--014} prove exact
properties of canonical Hermitian $2\times2$ matrices and bound maximum sorted-band
error by the unnormalized entrywise Frobenius distance. Numerical eigensolvers and
meshes belong to the numerical-verification layer.

\section{Checked certificates}

The certificate layer contains exact data structures and checkers whose soundness is
proved in Lean.  The rank-two separation families
\decl{BG-CERT-THM-001--009} and \decl{BG-CERT-THM-013--016} connect Boolean
acceptance to exact witness, containment, and strict-separation propositions, including
uniform box conclusions under the encoded endpoint hypotheses.  The separate finite
Pareto family \decl{BG-CERT-THM-010--012} proves exact dominance, threshold-feasibility,
and efficient-mask checker equivalences for one supplied finite objective table; it does
not establish external inventory coverage.

No retained runtime certificate is created by these theorems. A future external package
must still freeze its schema, independently verify its producer, justify exact or
outward-rounded inputs, retain hashes and provenance, and bind those values to the exact
checker proposition.

\section{Interval dispositions}

Declarations \decl{BG-CERT-THM-017--025} use three outcomes for upper-bound intervals:
accepted, rejected, and unresolved.  Acceptance requires a valid interval whose upper
endpoint passes.  For valid endpoints, rejection occurs when the lower endpoint lies
strictly above the budget.  Invalid endpoint order also fails closed to \emph{rejected},
but no substantive failure conclusion follows without an explicit containment premise;
a malformed interval cannot contain an exact value.  A valid interval overlapping the
threshold is unresolved.

Joint acceptance requires every objective to accept for the same candidate. One
rejected objective supplies joint rejection; otherwise at least one unresolved objective
leaves the joint result unresolved. Soundness requires explicit containment of each exact
value. Canonical serialization and binary64 storage require a separate enclosure argument
to provide that containment.

\section{Finite first-passing semantics}

Declarations \decl{BG-DECISION-THM-014--016} define and characterize first passage on a
complete supplied finite inventory. Every earlier declared cutoff must fail, and one
common cutoff must satisfy all objectives in the joint version. No monotonicity premise
is needed.

This result governs a complete declared finite radius, support, or cutoff inventory.
Continuum and asymptotic crossover claims require additional analytical results.

\section{Finite defect and loop algebra}

The rank-one resolvent family \decl{QM-RESOLVENT-THM-001--014} proves exact outer-product,
determinant, Sherman--Morrison, secular-zero, kernel-vector, and Hermiticity results under
explicit invertibility and denominator premises. A supplied secular root has exact
consequences; root existence, uniqueness, monotonicity, and infinite-volume behavior are
not proved.

The finite gauge-loop and projector--Bargmann declarations
\decl{SS-GAUGE-THM-001--023} prove closed-product covariance, commutative invariance
where applicable, Wilson-loop determinant invariance, rank-one projector identities,
and the ordered projector-product/Bargmann-product crosswalk. Nonvanishing, Wilson-loop
eigenphase invariance, phase unwrapping, Chern integrality, continuum topology, and
Wannier obstruction are not established by those results.

\chapter{Mechanization Status and Open Obligations}
\label{ch:mechanization-status}

\section{Status source}

The authoritative declaration inventory is \modulepath{docs/theorem-manifest.md}. The
named axiom reports imported by \modulepath{ProjectKoios/Audit.lean} provide the audit
surface for proved theorems and lemmas. This chapter intentionally does not duplicate a
mutable declaration count. A later export must take its count, toolchain identity, and
module hashes from one exact Formal revision.

The current basis of this working companion is merge commit
\decl{b56d6b8e7005515220e6c9bf42851acccb023178}, tree
\decl{a94b113d3300ff8e692f295da82945f48fb7d6ed}.  At that tree,
\path{ProjectKoios/BandGauge/Fourier/GaugeTruncationCounterexample.lean} has SHA-256
\decl{3b4d9083c04592ff9982cbddf891c8243cc599ad6bc32f5e4a10f708bec6186d}, and
\path{ProjectKoios/BandGauge/Fourier/ResidualChannels.lean} has SHA-256
\decl{d2da775a99edf7dea2d8955a121d4929cc183be4fcf64a515ab5e002c019a28b}.
No claim about later revisions follows from these identities.

\section{Checked mathematical families}

At the current snapshot, checked reusable families include:

\begin{itemize}
  \item finite Fourier transform, reconstruction, Parseval, truncation, Hermitian
        pairing, an exact gauge--truncation counterexample, and deliberate-truncation
        power, root-mean-square, pointwise, and maximum tail results;
  \item finite frames, sewing, aligned differences, and alignment-family separation;
  \item weighted Frobenius, spectral, admissible-set, and rank-two loss geometry;
  \item exact rational, enclosure, box, Pareto, and interval decision certificates;
  \item finite first-passing decisions;
  \item finite folding, spin lifting, route commutativity, and tight-binding algebra;
  \item rank-one resolvent defects and rank-two spectral perturbation;
  \item finite gauge loops, projector products, and Bargmann identities; and
  \item finite quantum-state, measurement, and partial-operator foundations.
\end{itemize}

Each family retains its own assumptions and nonclaims. Membership in this list does not
mean that every prospective external statement in the same subject has been encoded.

\section{Near-term exact theorem opportunities}

Several bounded extensions would sharpen the proof program without importing runtime or
material authority.

\begin{enumerate}
  \item \openobligation{Prove induced Euclidean operator-norm invariance under unitary
        conjugation in the exact finite setting.}
  \item \openobligation{Package the gauge invariance of a conformant two-path residual
        for both Frobenius and operator norms.}
  \item \openobligation{Develop a general finite pointwise-gauge convolution and
        support formula beyond the proved two-point counterexample.}
  \item \openobligation{Package shell-resolved Frobenius invariance under one constant
        unitary that does not mix shell labels.}
  \item \openobligation{Develop exact finite Feshbach or Schur-complement elimination
        separately from spectral-gap and norm-bound certificates.}
  \item \openobligation{Formalize a bounded projection-induced support-growth result
        without claiming infinite-volume decay or topology.}
\end{enumerate}

\openobligation{These are proposed theorem tasks, not declarations already present in
the manifest.}

\section{Analytical and scientific obligations kept open}

The current finite library does not establish:

\begin{itemize}
  \item asymptotic locality or existence of a continuum crossover radius;
  \item discrete-to-continuum, infinite-volume, or homogenization convergence;
  \item existence of a physically adequate smooth or symmetry-compatible alignment;
  \item a runtime frame's exact orthonormality or a numerical interval's exact
        containment;
  \item Chern integrality, global branch unwrapping, or a Wannier obstruction;
  \item Quantum ESPRESSO or Wannier90 execution, convergence, or material adequacy;
  \item silicon-specific applicability, scientific validation, novelty, publication, or
        human acceptance.
\end{itemize}

A future result may close one item only through a separately reviewed theorem or evidence
contract. Closure of one item does not close the others.

\section{Export checklist}

Before extracting prose into an external chapter, record:

\begin{enumerate}
  \item the exact Formal repository, commit, tree, module paths, and hashes;
  \item every declaration ID used by the exported argument;
  \item the explicit assumptions and operand/sign conventions;
  \item the named axiom-audit result and any remaining trusted-base limitations;
  \item which surrounding claims remain prospective, numerical, scientific, or
        publication-dependent; and
  \item confirmation that the destination does not present the export as runtime or
        material validation.
\end{enumerate}

The destination may reorganize the exposition, but it must not silently strengthen the
formal statements or detach them from their immutable source identity.


\appendix
\chapter{Formal Declaration Crosswalk and Mechanization Status}
\label{app:formal-crosswalk}
\small
\sloppy
\hbadness=10000

\section{Purpose of the crosswalk}

This appendix binds the mathematical exposition to stable declaration identifiers,
exact Lean names, source modules, and statement-level proof dossiers. The identifiers
are editorial handles; the Lean declarations and their proof terms remain authoritative.
A range in the tables means that every identifier in the inclusive range is represented
in the theorem manifest. It does not authorize a stronger aggregate statement than the
individual declarations prove.

\section{Finite-frame sequence}

The definitions and ten theorems in Chapter~\ref{ch:finite-operator-foundations} are
implemented in
\modulepath{ProjectKoios/BandGauge/Alignment/FiniteFrame.lean}. Their statement dossiers
are stored under \modulepath{proofs/BandGauge/Alignment/}.

\begin{description}
  \item[\decl{BG-ALIGNMENT-THM-001}.]
    \path{ProjectKoios.BandGauge.Alignment.frameProjector_isHermitian};
    dossier \modulepath{BG-ALIGNMENT-THM-001-projector-hermitian.tex}.
    This is Theorem~\ref{thm:bg-alignment-001}; it has no orthonormality premise.

  \item[\decl{BG-ALIGNMENT-THM-002}.]
    \path{ProjectKoios.BandGauge.Alignment.frameProjector_mul_self};
    dossier \modulepath{BG-ALIGNMENT-THM-002-projector-idempotent.tex}.
    This is Theorem~\ref{thm:bg-alignment-002}; exact frame orthonormality is required.

  \item[\decl{BG-ALIGNMENT-THM-003}.]
    \path{ProjectKoios.BandGauge.Alignment.rightUnitaryFrame_isOrthonormal};
    dossier \modulepath{BG-ALIGNMENT-THM-003-right-unitary-orthonormal.tex}.
    This is Theorem~\ref{thm:bg-alignment-003}.

  \item[\decl{BG-ALIGNMENT-THM-004}.]
    \path{ProjectKoios.BandGauge.Alignment.frameProjector_rightUnitaryFrame};
    dossier \modulepath{BG-ALIGNMENT-THM-004-projector-right-unitary.tex}.
    This is Theorem~\ref{thm:bg-alignment-004}; it does not require orthonormality.

  \item[\decl{BG-ALIGNMENT-THM-005}.]
    \path{ProjectKoios.BandGauge.Alignment.compressOperator_rightUnitaryFrame};
    dossier \modulepath{BG-ALIGNMENT-THM-005-compression-covariance.tex}.
    This is Theorem~\ref{thm:bg-alignment-005}.

  \item[\decl{BG-ALIGNMENT-THM-006}.]
    \path{ProjectKoios.BandGauge.Alignment.compressOperator_isHermitian};
    dossier \modulepath{BG-ALIGNMENT-THM-006-compression-hermitian.tex}.
    This is Theorem~\ref{thm:bg-alignment-006}; frame orthonormality is not required.

  \item[\decl{BG-ALIGNMENT-THM-007}.]
    \path{ProjectKoios.BandGauge.Alignment.reconstructCompressedOperator_eq_projector_sandwich};
    dossier \modulepath{BG-ALIGNMENT-THM-007-reconstruction-projector-sandwich.tex}.
    This is Theorem~\ref{thm:bg-alignment-007}.

  \item[\decl{BG-ALIGNMENT-THM-008}.]
    \path{ProjectKoios.BandGauge.Alignment.reconstructCompressedOperator_rightUnitaryFrame};
    dossier \modulepath{BG-ALIGNMENT-THM-008-reconstruction-right-unitary.tex}.
    This is Theorem~\ref{thm:bg-alignment-008}; it does not require orthonormality.

  \item[\decl{BG-ALIGNMENT-THM-009}.]
    \path{ProjectKoios.BandGauge.Alignment.frameProjector_mul_frame};
    dossier \modulepath{BG-ALIGNMENT-THM-009-projector-fixes-frame.tex}.
    This is Theorem~\ref{thm:bg-alignment-009}.

  \item[\decl{BG-ALIGNMENT-THM-010}.]
    \path{ProjectKoios.BandGauge.Alignment.compressOperator_frameProjector};
    dossier \modulepath{BG-ALIGNMENT-THM-010-compressed-projector-identity.tex}.
    This is Theorem~\ref{thm:bg-alignment-010}.
\end{description}

The definitions used by the sequence are
\decl{IsOrthonormalFrame}, \decl{frameProjector}, \decl{rightUnitaryFrame},
\decl{compressOperator}, and \decl{reconstructCompressedOperator}. Exact
orthonormality is a hypothesis, not a conclusion about a runtime frame.

\section{Bloch alignment and alignment families}

\begin{table}[ht]
\centering
\small
\begin{tabularx}{\textwidth}{@{}Z Y Y@{}}
\toprule
Declaration range & Source module & Mathematical role \\
\midrule
\decl{BG-BLOCH-THM-012--020}
  & \modulepath{ProjectKoios/BandGauge/Bloch/AlignedDifference.lean}
  & pointwise difference identities, shared-sewing closure, sign reversal, recovery,
    zero criteria, and aligned-difference sewing \\
\decl{BG-ALIGNMENT-DEF-007--016}
  & \modulepath{ProjectKoios/BandGauge/Alignment/Family.lean}
  & singleton, identity, global-unitary, and complete pointwise alignment families and
    their feasibility predicates \\
\decl{BG-ALIGNMENT-THM-011--019}
  & \modulepath{ProjectKoios/BandGauge/Alignment/Family.lean}
  & family membership, forward inclusions, and feasibility transfer \\
\decl{BG-ALIGNMENT-THM-020--029}
  & \modulepath{ProjectKoios/BandGauge/Alignment/FamilySeparation.lean}
  & strict family separation and exact reverse-feasibility counterexamples \\
\bottomrule
\end{tabularx}
\caption{Aligned-difference and alignment-family implementation crosswalk.}
\label{tab:alignment-family-implementation-crosswalk}
\end{table}

The corresponding dossiers are under \modulepath{proofs/BandGauge/Bloch/} and
\modulepath{proofs/BandGauge/Alignment/}. The strictness declarations use supplied point
and unitary inequalities. Their constructed acceptance predicates establish logical
separation of decision spaces; they do not certify the outcome of a scientific
application.

\section{Three-channel residual calculus}

The source-independent finite residual algebra used in
Theorem~\ref{thm:three-channel-residual-calculus} is implemented in
\modulepath{ProjectKoios/BandGauge/Hilbert/ResidualChannels.lean}.

\begin{table}[ht]
\centering
\small
\begin{tabularx}{\textwidth}{@{}Z Y Y@{}}
\toprule
Declaration range & Formal objects & Mathematical role \\
\midrule
\decl{BG-HILBERT-DEF-013--016}
  & \decl{FiniteResidualField}, \decl{interpolationResidual},
    \decl{truncationResidual}, \decl{totalResidual}
  & the common finite field type and the three signed residual channels \\
\decl{BG-HILBERT-DEF-017--018}
  & \decl{finiteMaximumNorm}, \decl{finiteRootMeanSquareNorm}
  & the finite maximum and root-mean-square summaries \\
\decl{BG-HILBERT-THM-016--017}
  & maximum-as-supremum and explicit RMS square-sum formula
  & normalization and interpretation of the two aggregate norms \\
\decl{BG-HILBERT-THM-018--021}
  & exact decomposition and pointwise, maximum, and RMS bounds
  & the three-channel identity and its norm-triangle consequences \\
\bottomrule
\end{tabularx}
\caption{Finite three-channel residual implementation crosswalk.}
\label{tab:three-channel-implementation-crosswalk}
\end{table}

The four principal declarations are:

\begin{description}
  \item[\decl{BG-HILBERT-THM-018}.]
    \path{ProjectKoios.BandGauge.Hilbert.totalResidual_eq_add};
    dossier \modulepath{BG-HILBERT-THM-018-residual-decomposition.tex}.
  \item[\decl{BG-HILBERT-THM-019}.]
    \path{ProjectKoios.BandGauge.Hilbert.norm_totalResidual_le_add};
    dossier \modulepath{BG-HILBERT-THM-019-pointwise-triangle.tex}.
  \item[\decl{BG-HILBERT-THM-020}.]
    \path{ProjectKoios.BandGauge.Hilbert.finiteMaximumNorm_totalResidual_le_add};
    dossier \modulepath{BG-HILBERT-THM-020-maximum-triangle.tex}.
  \item[\decl{BG-HILBERT-THM-021}.]
    \path{ProjectKoios.BandGauge.Hilbert.finiteRootMeanSquareNorm_totalResidual_le_add};
    dossier \modulepath{BG-HILBERT-THM-021-rms-triangle.tex}.
\end{description}

The dossiers are stored under \modulepath{proofs/BandGauge/Hilbert/}.  The exact identity
requires one common intermediate field.  The aggregate theorems are inequalities and do
not assert additivity of reported scalar errors.  The formal RMS definition is totalized
on an empty inventory; the scientific formulas in the main text use the ordinary
nonempty finite-inventory interpretation.

\section{Fourier and gauge families}

\begin{table}[ht]
\centering
\small
\begin{tabularx}{\textwidth}{@{}Z Y Y@{}}
\toprule
Declaration range & Source module & Mathematical role \\
\midrule
\decl{BG-FOURIER-THM-033--047}
  & \modulepath{ProjectKoios/BandGauge/Fourier/UnitaryEquivariance.lean}
  & fixed-unitary preservation and commutation with the finite Fourier and truncation
    constructions \\
\decl{BG-FOURIER-THM-048--054}
  & \modulepath{ProjectKoios/BandGauge/Fourier/PointwiseGauge.lean}
  & mode shift, scalar-conjugation boundary, coefficient change, and pointwise
    Frobenius-power preservation \\
\decl{BG-FOURIER-THM-055--066}
  & \modulepath{ProjectKoios/BandGauge/Fourier/GaugeTruncationCounterexample.lean}
  & exact two-point rank-two gauge--truncation counterexample \\
\decl{BG-FOURIER-THM-067--075}
  & \modulepath{ProjectKoios/BandGauge/Fourier/ResidualChannels.lean}
  & deliberate-truncation residual, omitted power, root-mean-square identities, and
    pointwise and finite-maximum bounds \\
\bottomrule
\end{tabularx}
\caption{Fourier and gauge implementation crosswalk.}
\label{tab:fourier-gauge-implementation-crosswalk}
\end{table}

The key declarations for Theorem~\ref{thm:two-point-gauge-truncation} are:

\begin{description}
  \item[\decl{BG-FOURIER-THM-062}.]
    \path{ProjectKoios.BandGauge.Fourier.pointwiseSwapConjugation_odd_coefficient};
    dossier \modulepath{BG-FOURIER-THM-062-gauged-odd-coefficient.tex}.
  \item[\decl{BG-FOURIER-THM-064}.]
    \path{ProjectKoios.BandGauge.Fourier.pointwiseSwapConjugation_odd_coefficient_ne_zero};
    dossier \modulepath{BG-FOURIER-THM-064-gauged-odd-mode-nonzero.tex}.
  \item[\decl{BG-FOURIER-THM-066}.]
    \path{ProjectKoios.BandGauge.Fourier.pointwiseSwapConjugation_zeroModeFourierTruncation_ne};
    dossier \modulepath{BG-FOURIER-THM-066-gauge-truncation-noncommutation.tex}.
\end{description}

The declaration sequence for Theorem~\ref{thm:finite-fourier-tail-bounds} is:

\begin{description}
  \item[\decl{BG-FOURIER-THM-067}.]
    \path{ProjectKoios.BandGauge.Fourier.truncationResidual_fourierReconstruction_eq}.
  \item[\decl{BG-FOURIER-THM-068}.]
    \path{ProjectKoios.BandGauge.Fourier.finiteFourierPower_truncationResidual_eq_omitted}.
  \item[\decl{BG-FOURIER-THM-069}.]
    \path{ProjectKoios.BandGauge.Fourier.finiteMatrixFrobeniusPower_truncationResidual_eq_omitted}.
  \item[\decl{BG-FOURIER-THM-070}.]
    \path{ProjectKoios.BandGauge.Fourier.finiteRootMeanSquareNorm_truncationResidual_eq_sqrt_omitted}.
  \item[\decl{BG-FOURIER-THM-071}.]
    \path{ProjectKoios.BandGauge.Fourier.finiteRootMeanSquareNorm_truncationResidual_sq_eq_omitted}.
  \item[\decl{BG-FOURIER-THM-072}.]
    \path{ProjectKoios.BandGauge.Fourier.finiteMatrixFrobeniusRootMeanSquare_truncationResidual_eq_sqrt_omitted}.
  \item[\decl{BG-FOURIER-THM-073}.]
    \path{ProjectKoios.BandGauge.Fourier.finiteMatrixFrobeniusRootMeanSquare_truncationResidual_sq_eq_omitted}.
  \item[\decl{BG-FOURIER-THM-074}.]
    \path{ProjectKoios.BandGauge.Fourier.norm_truncationResidual_fourierReconstruction_le_omitted}.
  \item[\decl{BG-FOURIER-THM-075}.]
    \path{ProjectKoios.BandGauge.Fourier.finiteMaximumNorm_truncationResidual_le_omitted}.
\end{description}

The dossiers follow the stable filenames
\modulepath{proofs/BandGauge/Fourier/BG-FOURIER-THM-067-truncation-channel.tex}
through
\modulepath{proofs/BandGauge/Fourier/BG-FOURIER-THM-075-finite-maximum-tail.tex}.
The two root-mean-square families require the explicit nonempty-domain premise recorded
in their Lean types. The maximum-norm result uses the ambient norm and is not an
entrywise matrix-Frobenius specialization.

\section{Reduction, certificates, and finite decisions}

\begin{table}[ht]
\centering
\small
\begin{tabularx}{\textwidth}{@{}Z Y Y@{}}
\toprule
Declaration range & Principal source & Scope represented in Chapter~\ref{ch:reduction-decision} \\
\midrule
\decl{BG-THM-001--037}
  & \modulepath{ProjectKoios/BandGauge/Geometry/}
  & common witnesses, compact attainment, zero/positive separation converses, rank-two
    admissible regions, and exact set distances \\
\decl{BG-HILBERT-THM-001--015}
  & \modulepath{ProjectKoios/BandGauge/Hilbert/}
  & ball geometry, weighted Frobenius algebra, and spectral/operator residual bridges \\
\decl{BG-FOUNDATION-THM-008--014}
  & \modulepath{ProjectKoios/BandGauge/Foundations/RankTwoSpectralPerturbation.lean}
  & canonical Hermitian rank-two spectra and Frobenius perturbation bounds \\
\decl{BG-CERT-THM-001--009}, \decl{BG-CERT-THM-013--016}
  & \modulepath{ProjectKoios/BandGauge/Certificates/}
  & exact, enclosure, and box separation certificates \\
\decl{BG-CERT-THM-010--012}
  & \modulepath{ProjectKoios/BandGauge/Certificates/FinitePareto.lean}
  & exact finite dominance, feasibility, and efficient-mask checks \\
\decl{BG-CERT-THM-017--025}
  & \modulepath{ProjectKoios/BandGauge/Certificates/IntervalDecision.lean}
  & accepted, rejected, and unresolved interval decisions and their soundness \\
\decl{BG-DECISION-THM-014--016}
  & \modulepath{ProjectKoios/BandGauge/Decision/FiniteCutoff.lean}
  & first passing and first joint-passing semantics on a finite inventory \\
\decl{BG-DECISION-DEF-012--027}
  & \modulepath{ProjectKoios/BandGauge/Decision/FiniteCompatibility.lean}
  & finite model/candidate inventories, infimum and witness compatibility, mandatory
    channels, conformant embeddings, metric separation, and first-compatible selection \\
\decl{BG-DECISION-THM-017--038}
  & \modulepath{ProjectKoios/BandGauge/Decision/FiniteCompatibility.lean}
  & infimum/attainment bridges, common-candidate and common-alignment semantics,
    compatibility transfer, finite-set distance attainment and zero/positive converses,
    and finite first-compatible theorems \\
\decl{QM-RESOLVENT-THM-001--014}
  & \modulepath{ProjectKoios/QuantumMechanics/FiniteDimensional/RankOneResolvent.lean}
  & exact rank-one resolvent, determinant, secular, kernel, and Hermiticity identities \\
\decl{SS-GAUGE-THM-001--023}
  & \modulepath{ProjectKoios/SolidState/FiniteLattice/}
  & finite gauge loops, determinant invariance, projector products, and Bargmann
    identities \\
\decl{SS-ROUTE-THM-001--010}
  & \modulepath{ProjectKoios/SolidState/FiniteLattice/RouteCommutativity.lean}
  & finite pullback, spin, folding, and defect-extraction route algebra \\
\bottomrule
\end{tabularx}
\caption{Declaration families supporting the reduction and decision narrative.}
\label{tab:reduction-decision-implementation-crosswalk}
\end{table}
\FloatBarrier

These declarations now package the generic finite spectral--operator compatibility
logic, including alignment-infimum and stronger common-witness semantics, explicit
model embeddings, finite central admissible sets, attained separation under nonemptiness,
zero/positive distance converses, and first-compatible selection.  No declaration in the
pinned snapshot instantiates the schema with the complete ICMSEP model hierarchy,
application-specific losses, tolerances, supports, channel contracts, or retained
certificate.  Chapter~\ref{ch:mechanization-status} records that proposition-specific
assembly as open.

\section{Mechanization status at the pinned snapshot}

The theorem manifest is the authoritative declaration inventory. At the pinned snapshot,
the repository audit covers 645 proved theorem or lemma rows and 649 unique named axiom
reports. The repository-owned Lean sources contain no \texttt{sorry}, \texttt{admit}, or
repository-owned \texttt{axiom} declaration. Individual reports may include the standard
foundational dependencies \texttt{propext}, \texttt{Classical.choice}, and
\texttt{Quot.sound}; the exact report for each declaration is recorded in the theorem
manifest and imported by \modulepath{ProjectKoios/Audit.lean}.

Checked reusable families include finite Fourier transformation and truncation, finite
frames and alignment families, weighted Frobenius and admissible-set geometry, exact and
interval certificates, finite first-passing decisions, finite lattice folding and route
algebra, rank-one resolvents, rank-two perturbation, finite gauge loops and Bargmann
identities, and finite quantum-state, measurement, and partial-operator foundations.
Membership in this inventory means that the exact Lean proposition is proved. It does
not imply runtime conformance, continuum validity, physical applicability, novelty, or
publication status.

\chapter{Provenance, Reproducibility, and Authority Boundaries}
\label{app:provenance}

\section{Immutable theorem snapshot}

The last immutable Project Koios Formal revision underlying this uncommitted proof-tranche
candidate is merge commit

\begin{center}
\decl{1078380fb608269a8094c7e86776a1c0d826333a},
\end{center}

with Git tree

\begin{center}
\decl{0e9695296ab1275df87e2b4b0309e3ef8c4d24b2}.
\end{center}

The candidate additionally modifies the geometry and finite-compatibility modules.  Its
integration commit and tree do not yet exist and must replace this predecessor boundary
after a separately authorized checkpoint.  Modules central to the alignment, residual,
and compatibility narrative have the following SHA-256 identities in the candidate
working tree:

\begin{description}
  \item[Admissible-set geometry.]
    \modulepath{ProjectKoios/BandGauge/Geometry/AdmissibleSets.lean},
    SHA-256
    \decl{e830b1b940ced88bc86b169261673337126583ee09e9ad9fb30ff580e9160fd8}.
  \item[Finite compatibility schema.]
    \modulepath{ProjectKoios/BandGauge/Decision/FiniteCompatibility.lean},
    SHA-256
    \decl{45e737324df4718497566d57d72ff9bd575f37d59a5f6768a101034b4c232a0d}.
  \item[Aligned operator differences.]
    \modulepath{ProjectKoios/BandGauge/Bloch/AlignedDifference.lean},
    SHA-256
    \decl{050e2cc7736e583c87f65488059f587b259fe5748960822b7f7eaba2485ffab8}.
  \item[Alignment families.]
    \modulepath{ProjectKoios/BandGauge/Alignment/Family.lean},
    SHA-256
    \decl{27b1c817bf10778a6d9032e6d5c8d20601aa9fa12f3c70ab560cdbf7dc672aab}.
  \item[Alignment-family separation.]
    \modulepath{ProjectKoios/BandGauge/Alignment/FamilySeparation.lean},
    SHA-256
    \decl{8811fffaca7411e525f37e710ad6c5b3ea9ba9064cfc74c3435344e8858caa71}.
  \item[Three-channel residual algebra.]
    \modulepath{ProjectKoios/BandGauge/Hilbert/ResidualChannels.lean},
    SHA-256
    \decl{38d0604da95a88a85ecc818037fedda0fb73ee9a51055bffa0a6891d8a4ee150}.
  \item[Gauge--truncation counterexample.]
    \modulepath{ProjectKoios/BandGauge/Fourier/GaugeTruncationCounterexample.lean},
    SHA-256
    \decl{3b4d9083c04592ff9982cbddf891c8243cc599ad6bc32f5e4a10f708bec6186d}.
  \item[Residual channels and finite tails.]
    \modulepath{ProjectKoios/BandGauge/Fourier/ResidualChannels.lean},
    SHA-256
    \decl{d2da775a99edf7dea2d8955a121d4929cc183be4fcf64a515ab5e002c019a28b}.
\end{description}

The predecessor revision and candidate module hashes make the pre-commit boundary
explicit, but they do not yet constitute one immutable integrated source state.  This
candidate is not publication-ready until an authorized checkpoint supplies a commit and
tree and this paragraph is refreshed.  No statement about a later revision follows
merely because a declaration retains the same informal name.  The commit, tree, hashes,
theorem-manifest entries, and proof dossiers must be reviewed together.

\section{Motivating ICMSEP problem statement}

The central spectral--operator compatibility question is taken from the owner repository
\nolinkurl{https://github.com/eragasa/ksdft2effmass.git}.  The source was inspected
read-only on 2026-10-29 at commit

\begin{center}
\decl{4835e7bf9161729cc8f0537bcc13db083896a6ae},
\end{center}

with tree

\begin{center}
\decl{fd641fb8a4efe7a5ef4d5da9c49f126e1e03e355}.
\end{center}

The checked-out branch was
\decl{work/monograph-current-library-capabilities}.  The owner worktree was broadly
dirty, so no repository-wide clean-state claim is made.  The two inspected source files
below were nevertheless byte-identical to the blobs at the recorded commit:

\begin{description}
  \item[Submitted-format extended abstract.]
    \modulepath{docs/publications/conferences/ICMSEP2026/ICMSEP-extended-abstract_02.docx},
    Git blob \decl{d23a7ae3bb137fa6d31f25ace37ba7e25fdd51a7},
    SHA-256
    \decl{4b1a2302ab47f7401b0531d9a89c29dfb05cebfc93e51eab1b8b2c96c351fa13}.
  \item[Markdown companion source.]
    \modulepath{docs/publications/conferences/ICMSEP2026/ksdft2effmass.ICMEP2026.abstract.md},
    Git blob \decl{dfcf5d0ea84a1f81843bd778bea3c8760b66b85d},
    SHA-256
    \decl{d1cf59403dc085d30f17e2bc46db321d9f25840c6f91a762ef157a59bf450ade}.
\end{description}

The extended abstract asks whether spectral and operator criteria admit one common
tight-binding Hamiltonian within declared tolerances, repeats that test over a nested
model hierarchy, and uses admissible-set separation when no common endpoint exists.
That problem statement controls the narrative hierarchy of this companion.  It does not
control Lean theorem types or proof status, and its prospective first-principles and
validation statements are not promoted here into completed scientific results.

\section{Seven-paper error-channel source snapshot}
\label{sec:seven-paper-source-snapshot}

The programmatic account in
Section~\ref{sec:seven-paper-error-channel-development} and
Appendix~\ref{app:seven-paper-error-channel-crosswalk} was inspected read-only in the
same owner repository at commit

\begin{center}
\decl{542ea2c4f16baaf4f3e905c2d833cc95db105dfe},
\end{center}

with tree

\begin{center}
\decl{ae9390fcda3d9142fab83b615104b55b094f8086}.
\end{center}

The checked-out branch was \decl{work/periodic1d-rank1-control-v3}; the worktree and all
listed source paths were clean at inspection on 2026-10-29.  The source identities are:

\begin{description}
  \item[Program roadmap.]
    \modulepath{docs/publications/conferences/ICMSEP2026/paper-program-roadmap.md};
    blob \decl{4db8f5617e7a557378da54e37fffdf03e29e2161};
    SHA-256 \decl{f0cbf0ee15fbaa5f8df80d0fe2d3ffabf02de7cfa3cb34c230c143d3e794b5f5}.
  \item[Experiment-to-paper allocation.]
    \modulepath{docs/publications/papers/experiment-to-paper-allocation.md};
    blob \decl{28d7fe01f176328cdd52268bd8339afce37d202c};
    SHA-256 \decl{e51954b0c2210617235b3e18d57577fe2d76691c099b3cc9d47ef6774fd9dd39}.
  \item[\decl{Ragasa2026a} v3.]
    \modulepath{docs/publications/papers/Ragasa2026a/v3/manuscript.tex};
    blob \decl{ac4dd332e9fc92fcd24c77875cdabbfc6e7852e9};
    SHA-256 \decl{02c20cf087cc7272ba45bb271d617a081582514c618bcdb6fe1c39d389d88d74}.
  \item[\decl{Ragasa2026b} v2.]
    \modulepath{docs/publications/papers/Ragasa2026b/v2/manuscript.tex};
    blob \decl{6e9dfef21e0cb7c33c56f8e01265510dcb9d88e1};
    SHA-256 \decl{d8d5951d8d7ed75355b5d94f37fcbddc70c1466897a5770dbe9f877e60016daf}.
  \item[\decl{Ragasa2026b} diagnostic definitions.]
    \modulepath{docs/publications/papers/Ragasa2026b/v2/appendix_diagnostics.tex};
    blob \decl{e8b62942faea8962cdc8767b1f201962588f3cfa};
    SHA-256 \decl{262b14b192a5893bd99a71d9457664a2c2d498f16af6e545b48273b26ac9e76d}.
  \item[\decl{Ragasa2026c} v1.]
    \modulepath{docs/publications/papers/Ragasa2026c/v1/manuscript.tex};
    blob \decl{4f1bd997ca3296f64472c9788f2668de4d1767bc};
    SHA-256 \decl{baf453722d96cf5b0360a0b80bfa1580465a1db7fe91167583943db10701baf4}.
  \item[\decl{Ragasa2026d} v1.]
    \modulepath{docs/publications/papers/Ragasa2026d/v1/manuscript.tex};
    blob \decl{248cbb6f32e3d1842cfc6be698a2001d1dbd6ff2};
    SHA-256 \decl{bb4aed2614777870c8607c7e67b43dbb133a51c672ac14cec6acf85004a18363}.
  \item[\decl{Ragasa2026e} v1.]
    \modulepath{docs/publications/papers/Ragasa2026e/v1/manuscript.tex};
    blob \decl{c896ffdb4c5ba52f5dbb8185b0738ceaaab755bd};
    SHA-256 \decl{72b60ef8cd2ba5505edeb4acce31ac004721c8ee5f6cee570729e19e870f8b76}.
  \item[\decl{Ragasa2026f} v1.]
    \modulepath{docs/publications/papers/Ragasa2026f/v1/manuscript.tex};
    blob \decl{af1e156b467a3ac4bf39e8b9e5a86f6896a4160a};
    SHA-256 \decl{a11a78b1a5cba72efb4e56695f4b9f1543369e97da7fcc149bee51f53913839c}.
  \item[\decl{Ragasa2026g} v1.]
    \modulepath{docs/publications/papers/Ragasa2026g/v1/manuscript.tex};
    blob \decl{c2b26e248d30f082276a0a3ca41a2b42d951b9ad};
    SHA-256 \decl{c52cc2f73f32c28a95ea2003004d36668b6cfa4c48d103d013c890af1cf42316}.
\end{description}

These owner sources establish editorial scope, terminology, and lifecycle status for the
seven-paper program.  They do not determine Lean theorem types, prove that a numerical
channel was evaluated correctly, or transfer controlled-model evidence into a material
claim.

\section{Authority hierarchy}

Four Formal-owned surfaces have distinct roles:

\begin{enumerate}
  \item Lean source under \modulepath{ProjectKoios/} owns exact definitions,
        propositions, and proof terms.
  \item \modulepath{docs/theorem-manifest.md} owns stable declaration identifiers and
        repository status.
  \item Files under \modulepath{proofs/} provide statement-level explanations,
        assumptions, provenance, and limitations.
  \item This companion synthesizes those declarations into a continuous mathematical
        argument.
\end{enumerate}

External manuscripts, calculation plans, schemas, and research notes may motivate the
formal program. They do not determine theorem statements or proof status. If prose in an
external source and a Lean type differ, the Lean type controls the formal claim and the
editorial crosswalk must record the discrepancy.

\section{Independent status axes}

The following questions remain logically independent:

\begin{itemize}
  \item Is an exact proposition implemented and proved in Lean?
  \item Which axioms does the compiled declaration report?
  \item Has its human-readable proof dossier been reviewed?
  \item Has an external certificate schema been frozen and independently verified?
  \item Does retained numerical evidence satisfy that schema?
  \item Are the theorem's assumptions applicable to the intended scientific system?
  \item Has novelty, publication, or human acceptance been established?
\end{itemize}

A positive answer on one axis cannot be transferred to another without additional
evidence. In particular, a clean Lean build does not certify runtime arrays, and a
reproducible numerical calculation does not prove an exact theorem unless an explicit
checked-certificate bridge connects its retained data to the theorem's premises.

\section{Dependency and toolchain boundary}

The pinned Formal snapshot uses Lean \decl{4.35.0}. Mathlib is pinned to immutable commit
\decl{56a005855ea17e0b0958f73ec801ae140b0a1679}; the dependency inspection reports nine
resolved revisions as commit identities. The durable dependency record is
\modulepath{lake-manifest.json}. A reproducibility review must inspect that file and must
not replace its immutable revisions with floating branches or tags.

The minimum repository validation boundary consists of:

\begin{enumerate}
  \item a successful Lake build for the exact Lean input tree, toolchain, and manifest;
  \item immutable dependency-pin inspection;
  \item an audit for admissions and repository-owned axioms;
  \item coverage of proved declarations by named axiom reports; and
  \item confirmation that generated build state is ignored rather than committed.
\end{enumerate}

A LaTeX build has a different meaning. It establishes that the exposition is
syntactically renderable and permits visual inspection of the resulting pages. It does
not elaborate Lean, validate a scientific model, or authenticate numerical evidence.

\section{Runtime certificate boundary}

The current runtime checked-certificate count for this companion is zero. No binary64
JSON payload, runtime array, eigensolver output, mesh, or serialized result is silently
reinterpreted as exact Lean input.

A future proposition-specific bridge must provide, at minimum:

\begin{itemize}
  \item a frozen schema that rejects unknown fields and missing mandatory data;
  \item exact values or outward-rounded enclosures with a documented rounding policy;
  \item separate producer and independently reviewed verifier implementations;
  \item retained canonical bytes, hashes, source revision, toolchain identity, and
        provenance;
  \item a precise statement of the Lean proposition established when the checker
        accepts; and
  \item strict separation between verification of the certificate and scientific
        assessment of the modeled system.
\end{itemize}

Canonical serialization can authenticate a byte sequence. It cannot make binary64 data
exact, establish enclosure, prove mesh completeness, or certify physical applicability.

\section{Requirements for editorial extraction}

A shortened publication or external chapter derived from this companion must record:

\begin{enumerate}
  \item the exact Formal repository, commit, tree, relevant module paths, and hashes;
  \item every declaration ID used by the extracted argument;
  \item the explicit assumptions and operand, sign, normalization, and domain
        conventions;
  \item the named axiom-audit result and remaining trusted-base limitations;
  \item which surrounding claims remain prospective, numerical, analytical,
        scientific, or publication dependent; and
  \item confirmation that the destination does not present the formal export as runtime
        verification or material validation.
\end{enumerate}

Editorial extraction may reorganize and shorten the exposition. It must not silently
strengthen a theorem, remove a premise, merge distinct alignment families, treat an
unresolved interval as decisive, or detach a statement from its immutable source
identity.

\chapter{Notation and Source-Manuscript Bridges}
\label{app:notation-bridges}
\small
\sloppy
\hbadness=10000

\section{Inspected source identity}

The retained-space, Bloch-frame, and aligned-difference bridges in this appendix were
prepared from the clean \texttt{ksdft2effmass} worktree on branch
\decl{work/monograph-representation-carriers-and-derivation-graphs} at commit
\decl{e3d8a51b30e2072aa4bb116eb46389e342fc00fc}. The inspected sources were the research
monograph chapters

\begin{itemize}
  \item \modulepath{chapters/finite-dimensional-foundations.tex};
  \item \modulepath{chapters/projection-and-retained-spaces.tex};
  \item \modulepath{chapters/composite-band-gauge-alignment.tex}; and
  \item \modulepath{chapters/operator-comparison.tex}.
\end{itemize}

The ICMSEP compatibility bridge and the three-channel residual sign bridge use the
separate immutable owner-source snapshots recorded in Appendix~\ref{app:provenance}.
Every bridge is editorial. The inspected manuscripts supply notation and motivation;
the Formal declarations identified in Appendix~\ref{app:formal-crosswalk} determine
exact statement and proof status.

\section{ICMSEP compatibility notation}

The ICMSEP extended-abstract source pinned in Appendix~\ref{app:provenance} supplies the
overall model-adequacy problem.  Table~\ref{tab:icmsep-compatibility-bridge} records its
relation to the notation of Chapters~\ref{ch:proof-architecture} and
\ref{ch:reduction-decision}.

\begin{table}[ht]
\centering
\small
\begin{tabularx}{\textwidth}{@{}l Y Y@{}}
\toprule
Proof companion & ICMSEP source & Mathematical role \\
\midrule
$\mathfrak M_j$
  & $\mathfrak M_j$
  & candidate tight-binding model class in the declared hierarchy \\
$H_j(\theta)$
  & $\mathbf H(\boldsymbol\theta)$
  & one candidate Hamiltonian in that class \\
$\mathcal L_{E,j}(\theta)$
  & $\mathcal L_E(\boldsymbol\theta)$
  & spectral loss \\
$\mathcal L_{H,j}(C,\theta)$
  & $\mathcal L_H(\mathbf C,\boldsymbol\theta)$
  & aligned operator loss \\
$\mathcal U_j$
  & $\mathcal U$
  & admissible alignment family \\
$\mathfrak A_{E,j}$, $\mathfrak A_{H,j}$
  & $\mathfrak A_{E,j}$, $\mathfrak A_{H,j}$
  & spectral and operator admissible sets \\
$\delta_j^*$
  & $\delta_j^*$
  & metric separation of the two admissible sets \\
\bottomrule
\end{tabularx}
\caption{Compatibility-problem notation bridge to the ICMSEP extended abstract.}
\label{tab:icmsep-compatibility-bridge}
\end{table}

The source writes a minimum over $\mathbf C\in\mathcal U$.  The companion uses an
infimum until attainment has been supplied, and then recovers the minimum formulation.
It also distinguishes literal nonempty intersection from acceptance under a positive
map tolerance.  These are scope clarifications, not changes to the source's motivating
question.

\section{Finite-dimensional retained-space notation}

The source manuscript and Chapter~\ref{ch:finite-operator-foundations} use the same
representation-neutral objects with the notation shown in
Table~\ref{tab:finite-manuscript-notation-bridge}.

\begin{table}[ht]
\centering
\small
\begin{tabularx}{\textwidth}{@{}l Y Y@{}}
\toprule
Proof companion & Source manuscript & Common mathematical object \\
\midrule
$V$
  & $V$
  & $D\times M$ retained-frame embedding \\
$|v_a\rangle$
  & $|\phi_a\rangle$
  & column $a$ of the retained frame \\
$P_V=VV^\dagger$
  & $P$ or $\widehat P_{\mathcal S}$
  & ambient projector onto the retained subspace \\
$C_V(H)=V^\dagger HV$
  & $H_{\mathcal S}$
  & operator matrix in retained coordinates \\
$R_V(H)=VC_V(H)V^\dagger$
  & $PHP$ embedded in the parent space
  & reconstructed retained operator in ambient coordinates \\
$V_G=VG$
  & $V'=VG$
  & right-unitary change of retained frame \\
\bottomrule
\end{tabularx}
\caption{Representation-neutral notation bridge to the inspected source manuscript.}
\label{tab:finite-manuscript-notation-bridge}
\end{table}

The source propositions labeled \texttt{PRF-05.01} and \texttt{PRF-05.02} correspond,
respectively, to projector invariance under a retained unitary and covariance of
compression. In the Formal snapshot these are
\decl{BG-ALIGNMENT-THM-004} and \decl{BG-ALIGNMENT-THM-005}. The nearby projector and
reconstruction claims are supported by \decl{BG-ALIGNMENT-THM-001--010} under the exact
premises listed in Appendix~\ref{app:formal-crosswalk}.

\section{Bloch-frame and Fourier notation}

The source composite-band chapter writes its momentum-dependent frame as
\[
  F(\kappa)=\begin{bmatrix}u_1(\kappa)&\cdots&u_M(\kappa)\end{bmatrix}.
\]
At corresponding momentum coordinates, the notation of the proof companion is
\[
  V(k)=F(\kappa)
  =\sum_{a=1}^M |u_a(\kappa)\rangle\,{}_{\mathrm{ret}}\!\langle a|.
\]
Thus $V(k)$ and $F(\kappa)$ are two names for the map that inserts retained-coordinate
vectors into an ambient Bloch fiber.

\begin{table}[ht]
\centering
\small
\begin{tabularx}{\textwidth}{@{}l Y Y@{}}
\toprule
Proof companion & Source manuscript & Common mathematical object \\
\midrule
$V(k)$
  & $F(\kappa)$
  & momentum-dependent retained-frame embedding \\
$|v_a(k)\rangle$
  & $u_a(\kappa)$
  & column $a$ of the composite Bloch frame \\
$P_V(k)$
  & $P(\kappa)=F(\kappa)F(\kappa)^\dagger$
  & ambient projector field \\
$C_{V(k)}(H(k))$
  & $H_F(\kappa)$
  & retained operator field \\
$T_R[h]$
  & $H_U(R)$ or $\langle w_{m0}|\widehat H|w_{nR}\rangle$
  & real-space Fourier block of the retained operator \\
\bottomrule
\end{tabularx}
\caption{Bloch and Fourier notation bridge to the inspected source manuscript.}
\label{tab:bloch-manuscript-notation-bridge}
\end{table}

The companion uses $G(k)$ for an internal unitary change of one retained frame and
reserves $U(k)$ for an identification between reference and candidate retained spaces.
The source composite-band chapter uses $U(\kappa)$ for its internal gauge. These symbols
must not be identified solely because their matrices are square and unitary: an internal
frame gauge and a two-model identification have different domains, provenance, and
adequacy obligations.

\section{Aligned-difference sign bridge}

When the source manuscript's pristine operator is identified with
$h_{\mathrm{ref}}$, its doped operator with $h_{\mathrm{cand}}$, and its map $U$ runs
from reference to candidate coordinates, the source writes
\[
  \Delta H_{d\to b}=U^\dagger h_{\mathrm{cand}}U-h_{\mathrm{ref}}.
\]
The convention adopted in this companion is
\[
  \Delta_U=h_{\mathrm{ref}}-U^{-1}h_{\mathrm{cand}}U.
\]
For a unitary $U$ with the corresponding direction, this gives
\[
  \Delta_U=-\Delta H_{d\to b}.
\]
Norms are insensitive to the sign, but signed perturbations, impurity potentials, and
response coefficients are not. Any editorial extraction must preserve both the operand
order and the direction assigned to $U$.

\section{Three-channel residual sign bridge}

Section~\ref{sec:three-channel-residual-calculus} and the Lean declarations use
reference-minus-approximation conventions:
\[
  \Delta_{\mathrm{int}}^{\mathrm F}=H_{\mathrm{ref}}-H_{\mathrm{int}},\qquad
  \Delta_{\mathrm{trunc}}^{\mathrm F}=H_{\mathrm{int}}-H_{\mathrm{tr}},\qquad
  \Delta_{\mathrm{total}}^{\mathrm F}=H_{\mathrm{ref}}-H_{\mathrm{tr}}.
\]
The signed spectral residuals in the inspected \decl{Ragasa2026a} source use the opposite
order at every step:
\[
  \Delta_{\mathrm{int}}^{\mathrm A}=H_{\mathrm{int}}-H_{\mathrm{ref}},\qquad
  \Delta_{\mathrm{trunc}}^{\mathrm A}=H_{\mathrm{tr}}-H_{\mathrm{int}},\qquad
  \Delta_{\mathrm{total}}^{\mathrm A}=H_{\mathrm{tr}}-H_{\mathrm{ref}}.
\]
Consequently
\[
  \Delta_x^{\mathrm A}=-\Delta_x^{\mathrm F}
  \quad\text{for }x\in\{\mathrm{int},\mathrm{trunc},\mathrm{total}\}.
\]
Simultaneous sign reversal preserves the additive chain and every norm reported in this
companion.  It does not permit signed fields from the two conventions to be combined
without conversion.  The inspected \decl{Ragasa2026b} source presents its three
channels through a symmetric error functional $\mathcal E(X,Y)$; that scalar notation
suppresses operand orientation but does not create scalar additivity.

\section{Status of the bridge}

The source-manuscript entries \texttt{PRF-*} are prospective editorial propositions.
They do not own the exact statements or proof status. The bridge records where the
source notation and the formal declarations agree, while the Lean type and proof term
remain authoritative whenever wording or scope differs.

\chapter{Seven-Paper Error-Channel Crosswalk}
\label{app:seven-paper-error-channel-crosswalk}
\small
\sloppy
\hbadness=10000

\section{Purpose and authority}

This appendix records how the finite three-channel residual calculus is used, refined,
or placed behind earlier gates across the seven-paper research program.  The immutable
owner-source identities are listed in
Section~\ref{sec:seven-paper-source-snapshot}.  The crosswalk is editorial: it does not
make an external manuscript a Formal authority, convert development output into retained
evidence, or assert that a prospective paper has completed its calculations.

The reusable exact core is
\[
  H_{\mathrm{ref}}-H_{\mathrm{tr}}
  =
  \bigl(H_{\mathrm{ref}}-H_{\mathrm{int}}\bigr)
  +
  \bigl(H_{\mathrm{int}}-H_{\mathrm{tr}}\bigr).
\]
It applies when all three fields share one domain, codomain, alignment convention, and
intermediate object.  Its norm consequences are inequalities.  The signed
\decl{Ragasa2026a} residuals reverse all three operand orders relative to the Formal
convention; Appendix~\ref{app:notation-bridges} records the simultaneous sign bridge.
The deliberate-truncation Fourier theorems identify only the second summand with an
omitted finite Fourier tail.

\section{Paper-by-paper development}

\begin{table}[ht]
\centering
\footnotesize
\begin{tabularx}{\textwidth}{@{}l Y Y@{}}
\toprule
Program item & Error-channel role & Checked-core relationship and current boundary \\
\midrule
\decl{Ragasa2026a}
  & Establishes separate cyclic-conversion, interpolation, deliberate-truncation, and
    total consequences in the controlled rank-two band-equivalence/locality study.
  & \emph{Formal relation:} the interpolation--truncation--total subchain instantiates
    the exact residual calculus when its fields are conformant; cyclic conversion
    remains additional, and only deliberate truncation has the Fourier-tail bridge.
    \emph{Boundary:} analytic results and bounded development diagnostics exist, but no
    retained active rank-two run or evidence package exists. \\
\decl{Ragasa2026b}
  & Carries explicit three-channel decompositions across orthogonal and nonorthogonal
    two- and three-dimensional controls, complete physical shells, and synthetic
    strain.
  & \emph{Formal relation:} the same three-field algebra separates interpolation floors
    from deliberate support restriction; dimension and geometry do not change the
    identity, but domains and supports remain application data.  \emph{Boundary:}
    verified development ladders are replaceable development output, not retained
    production evidence or material validation. \\
\decl{Ragasa2026c}
  & Refines the ledger into projection, disentanglement, downfolding, frame or gauge
    restriction, conversion, interpolation, truncation, and total consequences.
  & \emph{Formal relation:} the three-channel calculus is a nested final chain and does
    not certify upstream subspace or alignment operations.  \emph{Boundary:} a rank-two
    successor protocol and rank-$N$ development code and plan exist, but there is no
    stable result format, independent full-result check, retained run, or evidence
    package. \\
\bottomrule
\end{tabularx}
\caption{Controlled-model development of the residual-channel program.}
\label{tab:seven-paper-error-channel-crosswalk-controls}
\end{table}
\FloatBarrier

\begin{table}[ht]
\centering
\footnotesize
\begin{tabularx}{\textwidth}{@{}l Y Y@{}}
\toprule
Program item & Error-channel role & Checked-core relationship and current boundary \\
\midrule
\decl{Ragasa2026d}
  & Organizes the proposed bulk-silicon decision around parent-model, numerical or
    discretization, local-extraction, and model-reduction error.
  & \emph{Formal relation:} a three-channel chain may support Wannier interpolation and
    support truncation but is not the four-channel application budget.  \emph{Boundary:}
    pre-results; the tutorial preview is not production evidence, and no accepted
    production parent or proposition-specific Formal certificate exists. \\
\decl{Ragasa2026e}
  & Places geometry, projection, localization, common-space alignment, energy
    referencing, Fourier transformation, finite size, and later truncation in separate
    channels for pristine--defect subtraction.
  & \emph{Formal relation:} aligned and three-stage residuals become meaningful only
    after a qualified common retained space and scalar reference are supplied.
    \emph{Boundary:} pre-results; no common-space defect experiment package or retained
    defect operator exists. \\
\decl{Ragasa2026f}
  & Studies topological and symmetry obstructions to constructing the global alignment
    required by later operator comparisons.
  & \emph{Formal relation:} this is a premise gate; without the required map, the
    intended global aligned three-channel residual is not yet defined.  \emph{Boundary:}
    pre-results; no selected obstruction family, invariant protocol, or retained
    classification exists, and current Formal gauge-loop results do not establish
    continuum topology. \\
\decl{Ragasa2026g}
  & Proposes nested SOC impurity-model comparisons with parent, discretization,
    alignment, finite-supercell, impurity-reduction, continuum, operator, subspace, and
    observable errors kept visible.
  & \emph{Formal relation:} three-channel decompositions may be instantiated only after
    matched pristine and doped SOC branches, common-space alignment, and energy
    referencing pass their gates.  \emph{Boundary:} pre-results; no authorized matched
    pristine-Si/Si:P/Si:B SOC experiment or retained material result exists. \\
\bottomrule
\end{tabularx}
\caption{Application and obstruction boundaries of the residual-channel program.}
\label{tab:seven-paper-error-channel-crosswalk-applications}
\end{table}
\FloatBarrier

\section{Source locators}

The principal manuscript locators are:

\begin{itemize}
  \item \modulepath{docs/publications/papers/Ragasa2026a/v3/manuscript.tex};
  \item \modulepath{docs/publications/papers/Ragasa2026b/v2/manuscript.tex} and
        \modulepath{docs/publications/papers/Ragasa2026b/v2/appendix_diagnostics.tex};
  \item \modulepath{docs/publications/papers/Ragasa2026c/v1/manuscript.tex};
  \item \modulepath{docs/publications/papers/Ragasa2026d/v1/manuscript.tex};
  \item \modulepath{docs/publications/papers/Ragasa2026e/v1/manuscript.tex};
  \item \modulepath{docs/publications/papers/Ragasa2026f/v1/manuscript.tex}; and
  \item \modulepath{docs/publications/papers/Ragasa2026g/v1/manuscript.tex}.
\end{itemize}

Program order and lifecycle boundaries are controlled by
\modulepath{docs/publications/conferences/ICMSEP2026/paper-program-roadmap.md} and
\modulepath{docs/publications/papers/experiment-to-paper-allocation.md} at the source
snapshot pinned in Appendix~\ref{app:provenance}.  Individual manuscript prose is not
used to override the allocation ledger's distinction among analytic results,
development diagnostics, retained evidence, prospective work, and pre-results status.

\section{Interpretation rules}

Four rules govern use of this crosswalk.

\begin{enumerate}
  \item \textbf{One conformant chain.}  The exact identity requires one common
        reference, intermediate field, truncated field, evaluation inventory, sign
        convention, and any required alignment.  Separately optimized channels cannot
        be assembled after the fact.
  \item \textbf{Vector identity, scalar bounds.}  Residual fields add exactly; their
        norms generally do not.  Maximum and RMS summaries obey triangle inequalities,
        and a small total channel may conceal cancellation between larger components.
  \item \textbf{Fourier-tail scope.}  Omitted-coefficient identities and bounds apply
        to deliberate truncation of the declared finite interpolant.  They do not
        certify cyclic conversion, interpolation, projection, disentanglement,
        alignment, finite-size, parent-model, or continuum errors.
  \item \textbf{No status transfer.}  A checked generic theorem does not authenticate a
        runtime field, and a development plot does not instantiate the theorem's exact
        hypotheses.  Papers D--G remain prospective or pre-results at the pinned owner
        snapshot.
\end{enumerate}

Accordingly, the three-channel calculus is best understood as a transversal
verification spine.  It supports operator admissibility wherever a conformant
three-stage approximation chain exists, while the common-Hamiltonian intersection
remains the outer adequacy proposition and each paper retains its scientifically
necessary broader error ledger.


\end{document}
